\subsection{有理性的判定准则}

对于 K 的每个子域 L，设 \(\operatorname{End}_{\mathrm{L}}(\mathrm{K})\) 表示把 K 视为 L 上的左向量空间时的自同态环；若 L 包含 \(K'\) ，则 \(\operatorname{End}_{\mathrm{L}}(\mathrm{K})\) 是 \(\operatorname{End}_{\mathrm{K}'}(\mathrm{K})\) 的一个子环。对于 \(\operatorname{End}_{\mathrm{K}'}(\mathrm{K})\) 的每个子集 M，存在 K 的一个包含 \(K'\) 且使得 M 含于 \(\operatorname{End}_{\mathrm{L}}(\mathrm{K})\) 的最大子域 L，即由这样的 \(\xi \in K\) 组成的集合：\(\phi(\xi\eta) = \xi.\phi(\eta)\) 对一切 \(\eta \in K\) 与一切 \(\phi \in M\) 成立（立即可验证此集合是一个环；另一方面，把前述关系中的 \(\eta\) 换成 \(\xi^{-1}\eta\) ，即得 \(\phi(\xi^{-1}\eta) = \xi^{-1}.\phi(\eta)\) （当 \(\xi \neq 0\) ）。我们称这个域为 M 在 K 中的中心化子，并记作 \(\chi(\mathcal{M})\) 。现在设 V 是带 \(K'\) -结构 \(V'\) 的右向量 K-空间。对一切 \(\phi \in \operatorname{End}_{\mathrm{K}'}(\mathrm{K})\) ，存在且仅存在 Z-模 V 的一个自同态 \(\phi_{V}\) ，使得 \(\phi_{\mathbf{V}}(x'.\xi) = x'.\phi(\xi)\) 对 \(x' \in \mathbf{V}'\) 与 \(\xi \in \mathbf{K}\) 成立：因为在第 1 号中曾定义过一个 Z-同构 \(\lambda\) ，它定义在 \(\mathbf{V}' \otimes_{\mathbb{K}'} \mathbf{K}\) 上且映到 V 上，把 \(x' \otimes \xi\) 映为 \(x'\) . \(\xi\) ，而 \(\phi_{\mathbf{V}}\) 必然等于 \(\lambda \circ (1_{\mathbf{V}'} \otimes \phi) \circ \lambda^{-1}\) 。

定理 1. 设 \(\mathcal{M}\) 是 \(\operatorname{End}_{\mathbb{K}'}(\mathbf{K})\) 的一个子集，而 \(\mathbf{L} = \chi(\mathcal{M})\) 是 \(\mathbf{K}\) 中作为 \(\mathcal{M}\) 的中心化子的那个子域。

\begin{enumerate}
\def\labelenumi{(\roman{enumi})}
\item
  设 V 是带 K'-结构的右向量 K-空间。向量 \(x \in V\) 在 L 上有理，当且仅当 \(\phi_{\mathbf{v}}(x.\eta) = x.\phi(\eta)\) 对一切 \(\phi \in M\) 与一切 \(\eta \in K\) 成立。V 的向量子 K-空间 W 在 L 上有理，当且仅当 \(\phi_{\mathbf{v}}(\mathbf{W}) \subset \mathbf{W}\) 对一切 \(\phi \in M\) 成立。
\item
  设 \(V_{1}\) , \(V_{2}\) 是两个各带 \(K^{\prime}\) -结构的右向量 K-空间。K-线性映射 f 从 \(V_{1}\) 到 \(V_{2}\) 在 L 上有理，当且仅当 \(f(\phi_{\mathrm{V}_{1}}(x_{1})) = \phi_{\mathrm{V}_{2}}(f(x_{1}))\) 对一切 \(x_{1} \in V_{1}\) 与一切 \(\phi \in M\) 成立。
\end{enumerate}

我们首先对 \(x\) 证明断言 (i)。设 \(B\) 是 \(V\) 的在 \(K'\) 上有理的一组基，并记 \(x = \sum_{b \in B} b \cdot \xi_b\) ；那么，对于 \(\phi \in \mathcal{M}\) 和 \(\eta \in K\) ，

\[
\phi_ {V} (x. \eta) - x. \phi (\eta) = \sum_ {b \in B} b. (\phi (\xi_ {b} \eta) - \xi_ {b}. \phi (\eta))
\]

因此，关系

\[
" \text { for   all } \phi \in \mathcal {M} \text { and   all } \eta \in K, \phi_ {V} (x. \eta) = x. \phi (\eta)"
\]

和

\[
" \text { for   all } \phi \in \mathcal {M}, \text { all } b \in \mathrm{B} \text { and   all } \eta \in \mathrm{K}, \phi (\xi_ {b} \eta) = \xi_ {b}. \phi (\eta)"
\]

是等价的。这些关系中的第二个意味着对所有 \(b \in \mathbf{B}\) , \(\xi_b \in \chi(\mathcal{M})\) ，这证明了(i)的第一个断言。

接下来证明 (ii)。\(f\) 在 \(L\) 上有理，当且仅当对所有 \(x_1' \in V_1\) 有理于 \(K', f(x_1')\) 是 \(V_2\) 中一个在 \(L\) 上有理的向量；这将蕴含 \(f(x_1)\) 在 \(L\) 上有理，只要 \(x_1\) 是 \(V_1\) 中在 \(L\) 上有理的向量，这样的向量乃是以 \(L\) 中系数对 \(K'\) 上有理的向量作的线性组合。由论证的第一部分，上述条件等价于关系式

\[
f (x _ {1} ^ {\prime}) \cdot \phi (\eta) = \phi_ {\mathrm{v} _ {2}} (f (x _ {1} ^ {\prime}) \cdot \eta) \quad \text { for } \phi \in \mathcal {M} \text { and } \eta \in \mathrm{K}\tag{4}
\]

也可以写作

\[
f (\phi_ {\mathrm{v} _ {1}} (x _ {1} ^ {\prime}. \eta)) = \phi_ {\mathrm{v} _ {2}} (f (x _ {1} ^ {\prime}. \eta)) \quad \text { for } \phi \in \mathcal {M} \text { and } \eta \in \mathrm{K}.\tag{5}
\]

由于 \(V_{1}\) 的每个元素都是以 K 中系数对 \(V_{1}\) 中在 \(K'\) 上有理的元素作的线性组合，条件 (5) 等价于

\[
f (\phi_ {\mathrm{v} _ {1}} (x _ {1})) = \phi_ {\mathrm{v} _ {2}} (f (x _ {1}))
\]

对所有 \(x_{1} \in V_{1}\) 与所有 \(\phi \in \mathcal{M}\) 。

最后，为证明 (i) 中的第二个断言，我们使用第 6 号引理 1：W 是某个 K-线性映射 \(g \colon \mathbb{W}_1 \to \mathbb{W}_2\) 的图像，而 W 在 L 上有理当且仅当映射 \(g\) 在 L 上有理（第 3 号，命题 4）。由 (ii)，\(g\) 在 L 上有理，当且仅当 \(g(\phi_{\mathbb{W}_1}(x_1)) = \phi_{\mathbb{W}_2}(g(x_1))\) 对所有 \(x_1 \in \mathbb{W}_1\) 与所有 \(\phi \in \mathcal{M}\) 成立；由于 \(\phi_{\mathbb{V}} = \phi_{\mathbb{W}_1} \times \phi_{\mathbb{W}_2}\) ，上述条件意味着 \(g\) 的图像 W 在 \(\phi_{\mathbb{V}}\) （对所有 \(\phi \in \mathcal{M}\) ）之下稳定。

\section{仿射空间与射影空间}

\subsection{仿射空间的定义}

定义 1. 给定域 K 上的左（相应地，右）向量空间 T，相伴于 T 的仿射空间是加法群 T 的任意齐性空间 E（I，§5，第 5 号），使得 0 是 T 中唯一保持 E 的所有元素不变的算子（即 T 在 E 上忠实且传递地作用）。在这些条件下，T 称为 E 的平移空间，其元素称为 E 的平移（或 E 的自由向量）。

在下文中，我们将仅限于讨论 T 是 K 上的左向量空间的情形。仿射空间 E 的平移向量空间 T 的（在 K 上的）维数称为 E 的（在 K 上的）维数，记为 dim E 或 dim \(_{K}\) E。一维（相应地，二维）仿射空间称为仿射直线（相应地，仿射平面）。仿射空间的元素也称为点。

在定义 1 的条件下，对于 \(\mathbf{t} \in \mathbf{T}\) 与 \(a \in \mathbf{E}\) ，我们将用 \(\mathbf{t} + a\) 或 \(a + \mathbf{t}\) 表示点 \(a\) 在 \(\mathbf{t}\) 下的像。则关系式

\[
\mathbf {s} + (\mathbf {t} + a) = (\mathbf {s} + \mathbf {t}) + a, \quad 0 + a = a\tag{1}
\]

对 \(\mathbf{s} \in \mathrm{T}\) , \(\mathbf{t} \in \mathrm{T}\) , \(a \in \mathrm{E}\) 成立。映射 \(x \mapsto x + \mathbf{t}\) 是 \(\mathbf{E}\) 到其自身上的一个双射，我们把它与 \(\mathbf{t}\) 等同。此外，定义 1 还蕴含：对所有 \(a \in \mathrm{E}\) ，映射 \(\mathbf{t} \mapsto \mathbf{t} + a\) 是 \(\mathrm{T}\) 到 \(\mathbf{E}\) 上的一个双射。换言之，给定两个点 \(a\) , \(b\) （属于 \(\mathbf{E}\) ），存在且仅存在一个平移 \(\mathbf{t}\) 使得 \(b = \mathbf{t} + a\) ；我们把这个平移记作 \(b - a\) ；则公式

\[
a - a = 0, \quad a - b = - (b - a), \quad b = (b - a) + a\tag{2}
\]

\[
(c - b) + (b - a) = c - a
\]

对 \(a \in \mathbf{E}\) , \(b \in \mathbf{E}\) , \(c \in \mathbf{E}\) 成立。若四个点 \(a, b, a', b'\) （属于 \(\mathbf{E}\) ）满足 \(b - a = b' - a'\) ，则公式

\[
b ^ {\prime} = \left(b ^ {\prime} - b\right) + (b - a) + a = \left(b ^ {\prime} - a ^ {\prime}\right) + \left(a ^ {\prime} - a\right) + a
\]

以及 T 中加法的交换性表明 \(b' - b = a' - a\) .

给定一点 \(a \in E\) ，映射 \(x \mapsto x - a\) 是 \(E\) 到 \(T\) 上的一个双射；当 \(E\) 在此映射下与 \(T\) 等同时，我们就说把 \(E\) 视为通过取 \(a\) 作为 \(E\) 中的原点而得到的向量空间。反之，每个向量空间 \(T\) 都典范地具有一个相伴于 \(T\) 的仿射空间结构，即对应于子群 \(\{0\}\) （属于 \(T\) ）的齐性空间结构（I，§ 5，第 6 号）。

注记。本号中的定义以及后面的某些结果可立即推广到如下情形：我们不考虑向量空间 T，而是考虑任意一个带算子的交换群 T。

\subsection{重心计算}

命题 1. 设 \((x_{t})_{t\in I}\) 是仿射空间 \(\mathbf{E}\) 中的一族点， \((\lambda_t)_{t\in I}\) 是 \(\mathbf{K}\) 中一族具有有限支集的元素，使得 \(\sum_{i\in I}\lambda_i = 1\) （相应地， \(\sum_{i\in I}\lambda_i = 0\) ）。若 \(a\) 是 \(\mathbf{E}\) 的任意一点，则由下式定义的点 \(x\in \mathbf{E}\)

\[
x - a = \sum_ {i \in I} \lambda_ {i} (x _ {i} - a)
\]

（相应地，自由向量 \(\sum_{i\in I}\lambda_i(x_i - a))\) 与所考虑的点无关。

若 \(a'\) 是 \(E\) 的另一点，则

\[
\sum_ {i} \lambda_ {i} (x _ {i} - a ^ {\prime}) = \sum_ {i} \lambda_ {i} ((x _ {i} - a) + (a - a ^ {\prime})) = \sum_ {i} \lambda_ {i} (x - a) + \left(\sum_ {i} \lambda_ {i}\right) (a - a ^ {\prime}).
\]

若 \(\sum_{i}\lambda_{i} = 1\) ，则 \(\sum_{i}\lambda_{i}(x_{i} - a) = (x - a) + (a - a') = x - a'\) ；若 \(\sum_{i}\lambda_{i} = 0\) ，则 \(\sum_{i}\lambda_{i}(x - a') = \sum_{i}\lambda_{i}(x - a)\) ；由此得到命题。

在命题1的条件下，由下式定义的点x

\[
x - a = \sum_ {i \in I} \lambda_ {i} (x _ {i} - a)
\]

\(\left(\text{resp. 自由向量} \sum_{t \in I} \lambda_t(x_t - a)\right)\) 记作 \(\sum_{t \in I} \lambda_t x_t\) 。

这样，特别地就重新得到第 1 号中引入的记号 \(b - a\) 。当 \(\sum_{i} \lambda_{i} = 1\) 时，点 \(x = \sum_{i} \lambda_{i} x_{i}\) 称为点 \(x_{i}\) 在给定质量 \(\lambda_{i}\) 下的重心。

给定 \(m\) 个点 \(a_1, \ldots, a_m\) 于 \(\mathbf{E}\) ，其数目 \(m\) 不是 \(\mathbf{K}\) (V, § 1) 的特征的倍数，点 \(g = \sum_{i=1}^{m} \frac{1}{m} a_i\) 称为（滥用说法）点 \(a_i\) ( \(1 \leqslant i \leqslant m\) ) 的重心（对于 \(m = 2\) ，我们说“中点”而不说“重心”）；它由关系式

\[
\sum_ {i = 1} ^ {m} \left(a _ {i} - g\right) = 0.
\]

\subsection{仿射线性簇}

定义 2. 给定仿射空间 E，E 的子集 V 称为仿射线性簇（或简称线性簇，或 E 的仿射子集），如果对 V 的点的每一族 \((x_{t})_{t\in I}\) 以及 K 中具有有限支集的每一族元素 \((\lambda_t)_{t\in I}\) ，只要 \(\sum_{t\in I}\lambda_t = 1\) ，重心 \(\sum_{t\in I}\lambda_t x_t\) 就属于 V。

这等同于说，定义2的条件对V的每一有限点族都成立。

空集是一个线性簇；线性簇的任意交集仍是一个线性簇。

设 V 是 E 的一个非空子集， \(a\) 是 V 的一个点；关系

\[
x - a = \sum_ {i = 1} ^ {n} \lambda_ {i} (x _ {i} - a)
\]

意味着 \(x\) 是族的一个重心 \(\sum_{i=1}^{n} \lambda_i x_i + \left(1 - \sum_{i=1}^{n} \lambda_i\right)a\) ，该族由诸 \(x_i\) 与 \(a\) 组成。因此：

命题 2. 对于仿射空间 E 的非空子集 V，要使 V 成为线性簇，当且仅当在取 V 中一点为原点所得到的 E 的向量空间结构下，V 是向量子空间。

特别地，向量空间 T（视为仿射空间）的非空仿射线性簇恰好是 T 的向量子空间的平移；因此，T 的向量子空间就是包含 0 的线性簇。

设 V 是仿射空间 E 的非空线性簇；当 x 和 y 遍历 V 时，自由向量 x - y 的集合是 E 的平移空间 T 的一个向量子空间 D，称为 V 的方向：因为，若 \(a \in V\) ，则

\[
x - y = (x - a) - (y - a)
\]

于是我们的断言由命题2得出。显然，D在V上忠实且传递地作用，因此V典范地具有与D相联系的仿射空间结构。所谓仿射簇V的维数，我们指的是V在此仿射空间结构下的维数，即向量空间D的维数。维数为0的线性簇是E的点；维数为1（相应地，2）的线性簇称为E的直线（相应地，平面）。

每个向量 \(\neq0\) 属于某直线的方向，称为该直线的一个方向向量；其关于T的一组基的分量构成所谓该直线的方向参数组。

线性簇V在E中的余维数是指其方向D在T中的余维数；E中余维数为1的线性簇称为E的一个（仿射）超平面。

具有相同方向的两个线性簇称为平行的；这等同于说其中一个是由另一个通过平移得到的。若V是T（视为仿射空间）中的线性簇，则其方向是与V平行且包含0的线性簇。

命题 3. 给定一族点 \((a_{t})_{t\in I}\) （属于一个仿射空间 \(\mathbf{E}\) ），集合 \(\mathbf{V}\) ——由重心 \(\sum_{i\in I}\lambda_i a_i\) (( \(\lambda_i\) ) 具有有限支集， \(\sum_{i\in I}\lambda_i = 1\) ) 组成—— 是 \(\mathbf{E}\) 的一个线性簇。

若族 \((a_{t})\) 为空，则 \(\mathbf{V} = \varnothing\) ，这是由条件 \(\sum_{i}\lambda_{i} = 1\) 所致。因此可设族 \((a_{t})\) 非空，此时在 E 中取某个 \(a_{t}\) 作为原点，命题便是显然的。

线性簇 V 显然是包含诸 \(a_{t}\) 的最小线性簇；称它为由族 \((a_{t})\) 生成的线性簇，并称该族为 V 的一个生成系。

在命题 3 的记号下，假设族 \((a_{\iota})\) 非空，要使每个点 \(x \in V\) 的形如 \(x = \sum_{\iota} \lambda_{\iota} a_{\iota}\) 的表达式是唯一的，必要且充分条件是：用 \(\kappa\) 表示 I 的任意指标时，向量族 \(a_{\iota} - a_{\kappa}\) （其中 \(\iota\) 遍历指标集 \(\neq \kappa\) ）在 T 中是自由的。这时称 E 的点的族 \((a_{\iota})_{\iota \in I}\) 是仿射自由的（或称其元素构成一个仿射自由系，或称它们仿射无关），并称 \(\lambda_{\iota}\) 是 \(x\) 的指标为 \(\iota\) 的、关于仿射自由族 \((a_{\iota})\) 的重心坐标。

设 \((a_{i})_{i\in I}\) 是 E 中一族点，若它不是仿射自由的，则称它是仿射相关的。

命题 4. 要使点的非空族 \((a_{t})_{t\in I}\) 在仿射空间 \(E\) 中是仿射相关的，必要且充分条件是：存在一个族 \((\lambda_t)_{t\in I}\) ，其元素在 \(K\) 中且不全为零，具有有限支集，使得 \(\sum_{i\in I}\lambda_i = 0\) 且 \(\sum_{i\in I}\lambda_i a_i = 0\) 。

给定指标 \(\kappa \in I\) ，说向量族 \((a_{t} - a_{\kappa})\) （其中 \(t\) 遍历指标集 \(\neq \kappa\) ）在 \(T\) 中相关，意味着存在一个不全为零的标量族 \((\lambda_{t})_{t \neq \kappa}\) ，使得 \(\sum_{t \neq \kappa} \lambda_{t}(a_{t} - a_{\kappa}) = 0\) ，它也可写成 \(\sum_{t \in I} \lambda_{t} a_{t} = 0\) ，其中 \(\lambda_{\kappa} = -\sum_{t \neq \kappa} \lambda_{t}\) ，换言之 \(\sum_{t \in I} \lambda_{t} = 0\) 。

命题 5. 要使点的非空族 \((a_{t})_{t\in I}\) 在仿射空间 \(E\) 中是仿射自由的，必要且充分条件是：对每个指标 \(\kappa\in I\) ， \(a_{\kappa}\) 不属于由诸 \(a_{t}\) （指标 \(\neq\kappa\) ）生成的线性簇。

当 I 只有一个元素时，命题显然成立。否则，在 E 中取某个 \(a_{t}\) （指标 \(\neq\kappa\) ）作为原点，命题即由 §7 第 1 号的注记得出。

\subsection{仿射线性映射}

定义 3. 给定附属于同一域 K 上两个向量空间 T、T' 的两个仿射空间 E、E'，E 到 E' 的映射 u 称为仿射线性映射（或仿射映射），如果对 E 的点的每一族 \((x_{t})_{t\in I}\) 以及每一族 \((\lambda_{t})_{t\in I}\) （满足 \(\sum_{i\in I}\lambda_{i}=1\) ），

\[
u \left(\sum_ {i \in I} \lambda_ {i} x _ {i}\right) = \sum_ {i \in I} \lambda_ {i} u (x _ {i}).\tag{3}
\]

命题 6. 设 \(u\) 是 \(E\) 到 \(E'\) 的一个仿射映射。存在一个且仅一个线性映射 \(v\) ，它把 \(T\) 映到 \(T'\) ，使得

\[
u (x + \mathbf {t}) = u (x) + v (\mathbf {t})
\]

对所有 \(x \in \mathbf{E}\) ， \(\mathbf{t} \in \mathbf{T}\) 成立。

设 \(a\) 是 E 中任意一点。映射

\[
\mathbf {t} \mapsto u (a + \mathbf {t}) - u (a)
\]

是 T 到 T' 的一个线性映射 \(v_{a}\) ，因为我们可以写出

\[
\begin{array}{c} a + \lambda \mathbf {t} = \lambda (a + \mathbf {t}) + (1 - \lambda) a \\ a + \mathbf {s} + \mathbf {t} = (a + \mathbf {s}) + (a + \mathbf {t}) - a \end{array}
\]

且由 (3) 可得 \(v_{a}(\lambda \mathbf{t}) = \lambda v_{a}(\mathbf{t})\) 与 \(v_{a}(\mathbf{s} + \mathbf{t}) = v_{a}(\mathbf{s}) + v_{a}(\mathbf{t})\) 。此外，若 \(b\) 是 \(\mathbf{E}\) 中另一点，则 \(v_{a} = v_{b}\) ；因为关系

\[
(a + \mathbf {t}) - a + b = b + \mathbf {t}
\]

蕴含

\[
u (a + \mathbf {t}) - u (a) + u (b) = u (b + \mathbf {t})
\]

也就是 \(u(a + \mathbf{t}) - u(a) = u(b + \mathbf{t}) - u(b)\) 。由此得到 \(v\) 的存在性；唯一性是直接的。

v 称为与 u 相伴的、从 T 到 T' 的线性映射。反之，对 T 到 T' 的每个线性映射 v 以及每对点 \(a \in E\) , \(a' \in E'\) ，立即可验证

\[
x \mapsto a ^ {\prime} + v (x - a)
\]

是从 E 到 \(E'\) 的一个仿射映射，其相伴线性映射为 v。因此，说 u 是 E 到 \(E'\) 的仿射映射，也就意味着：若取 E 中任意一点 a 与点 \(u(a)\) （在 \(E'\) 中）作为原点，则 u 是关于由此所得两个向量空间的一个线性映射。

设 \(\mathbf{E}''\) 是第三个仿射空间， \(\mathbf{T}''\) 是它的平移空间， \(u'\) 是 \(\mathbf{E}'\) 到 \(\mathbf{E}''\) 的一个仿射映射，而 \(v'\) 是一个线性映射（从 \(\mathbf{T}'\) 到 \(\mathbf{T}''\) ），它与 \(u'\) 相伴。显然 \(u' \circ u\) 是 \(\mathbf{E}\) 到 \(\mathbf{E}''\) 的一个仿射映射；此外，对于 \(a \in \mathbf{E}\) 和 \(\mathbf{t} \in \mathbf{T}\) ，

\[
u ^ {\prime} (u (a + \mathbf {t})) = u ^ {\prime} (u (a) + v (\mathbf {t})) = u ^ {\prime} (u (a)) + v ^ {\prime} (v (\mathbf {t}))
\]

因此 \(v'\circ v\) 是 T 到 \(T''\) 的、与 \(u'\circ u\) 相伴的线性映射。仿射映射 u 是双射，必要且充分条件是它的相伴线性映射 v 是双射，且此时 \(u^{-1}\) 是一个仿射映射，其相伴线性映射为 \(v^{-1}\) 。

特别地，E 到自身的仿射双射构成一个群 G，称为 E 的仿射群。将 \(u \in G\) 映射到与 u 相伴的线性映射 v 的映射，由上述可知是 G 到线性群 \(\mathbf{GL}(\mathrm{T})\) 上的同态。若 u 是平移，则 v 是恒等映射，反之亦然。因此，上述同态的核是 E 的平移群 T，从而 T 是 G 的正规子群。

若 \(u \in G\) ，T 的自同构 \(t \mapsto utu^{-1}\) （其中 t 恒同于平移 \(x \mapsto x + t\) ）就是与 u 相伴的线性映射 v。对于 \(x \in E\) 和 \(t \in T\) ，根据定义

\[
x + u \mathbf {t} u ^ {- 1} = u (u ^ {- 1} (x) + \mathbf {t}) = u (u ^ {- 1} (x)) + v (\mathbf {t}) = x + v (\mathbf {t})
\]

因此 \(utu^{-1} = v(t)\) 。

设 \(a \in E\) ， \(G_a\) 表示 \(G\) 中由 \(u \in G\) （满足 \(u(a) = a\) ）组成的子群。若把 \(E\) 与 \(T\) 等同起来（取 \(a\) 为原点），则 \(G_a\) 就与 \(\mathbf{GL}(T)\) 等同。每个 \(u \in G\) 都可以唯一地表示成 \(u = t_1 u_1\) （相应地 \(u = u_2 t_2\) ）的形式，其中 \(u_1, u_2\) 属于 \(G_a\) ， \(t_1, t_2\) 属于 \(T\) ：因为令 \(t_1 = u(a) - a\) ，则 \(u^{-1} t_1 \in G_a\) ，由此得出 \(u_1\) 与 \(t_1\) 的存在性； \(u_2\) 与 \(t_2\) 的存在性同理可得。唯一性由 \(G_a \cap T\) 化为 \(G\) 的单位元这一事实得出。此外，

\[
\mathbf {t} _ {1} u _ {1} = u _ {1} (u _ {1} ^ {- 1} \mathbf {t} _ {1} u _ {1})
\]

由此得 \(u_{2}=u_{1}\) ， \(t_{2}=u_{1}^{-1}t_{1}u_{1}\) 。最后，与 u 和 \(u_{1}\) 相伴的线性映射相同，因此，若如上所述把 \(G_{a}\) 与 \(\mathbf{GL}(T)\) 等同起来，则 \(u_{1}\) 就是从 T 到自身的、与 u 相伴的线性映射。由此可见，G 是 \(G_{a}\) 与 T 的半直积（I，§6，第 1 号）。

设 E，E' 是 K 上的两个仿射空间。在 E（相应地 E'）到 E' 的仿射映射 u 下，E（相应地 E'）的线性簇的直像（相应地逆像）是 E'（相应地 E）的线性簇；u 的秩按定义是 \(u(\mathrm{E})\) 的维数；它等于与 u 相伴的线性映射的秩。若 V，V' 分别是 E，E' 中具有相同有限维数 m 的线性簇，则存在一个仿射映射 \(u\) ，把 \(\mathbf{E}\) 映到 \(\mathbf{E}'\) ，使得 \(u(\mathbf{V}) = \mathbf{V}'\) ：在 \(\mathbf{E}\) 与 \(\mathbf{E}'\) 中分别取 \(\mathbf{V}\) 与 \(\mathbf{V}'\) 的点作为原点，然后在 \(\mathbf{E}\) （相应地 \(\mathbf{E}'\) ）中取一个基，使其前 \(m\) 个向量构成 \(\mathbf{V}\) （相应地 \(\mathbf{V}'\) ）的基，该命题便立即可由 §1 第 11 号命题 17 的推论 3 得出。

由于域 K 典范地具有 K 上的（一维）左向量空间结构，它可以看作一个一维仿射空间。仿射空间 D（在 K 上）到仿射空间 K 的仿射映射也称为仿射线性函数（或仿射函数）。若在 E 中取一点 \(a\) 为原点，则 E 上的每个仿射函数都可以唯一地写成 \(x \mapsto \alpha + v(x)\) ，其中 \(\alpha \in \mathbf{K}\) ，而 \(v\) 是由此所得向量空间 E 上的一个线性形式；因此 E 上的仿射函数构成 K 上一个维数为 \(1 + \dim \mathbf{E}\) 的右向量空间。若 \(u\) 是 E 上一个非常数仿射函数，且 \(\lambda \in \mathbf{K}\) ，则由 \(x \in \mathbf{E}\) （满足方程 \(u(x) = \lambda\) ）组成的集合是一个超平面；反之，对每一个

E 中的超平面 H，存在 E 上的一个仿射函数 \(u_{0}\) ，使得 \(\mathbf{H} = \bar{u}_{0}^{-1}(0)\) ，且每个满足 \(\mathbf{H} = \bar{u}^{-1}(0)\) 的仿射函数 u 都具有形式 \(u_{0}\mu\) ，其中 \(\mu \in K\) （§7，第 5 号，命题 11）。若 u 是 E 上的仿射函数，则方程为 \(u(x) = \alpha\) 与 \(u(x) = \beta\) 的超平面相互平行。

\subsection{射影空间的定义}

定义 4. 给定域 K 上的一个左（相应地，右）向量空间 V，由 V 导出的左（相应地，右）射影空间，记作 P(V)，是指补集 V - \{0\} （即 \{0\} 在 V 中的补集）关于等价关系 Δ(V) 所作的商，其中 Δ(V) 为：“存在 K 中的 λ ≠ 0，使得 y = λx（相应地，y = xλ）”，该关系定义在 V - \{0\} 中的 x 与 y 之间。

当 \(\mathbf{V} = \mathrm{K}_s^{n + 1}\) 时，我们也用 \(\mathbf{P}_n(\mathbf{K})\) 代替 \(\mathbf{P}(\mathrm{K}_s^{n + 1})\) ，用 \(\Delta_n(\mathbf{K})\) 代替 \(\Delta (\mathbf{V})\) 。

定义 4 也可以这样表述： \(\mathbf{P}(\mathbf{V})\) 是 V 中（过 0 的）直线所成的集合，且原点已被去掉；因此 \(\mathbf{P}(\mathbf{V})\) 典范地等同于 V 中（过 0 的）直线的集合。射影空间的元素称为该空间的点。

当 V 的维数为 n 时，若 n 有限，整数 n - 1 称为射影空间 \(\mathbf{P}(\mathbf{V})\) 的维数，否则称为基数 n；这个基数记作 \(\dim_{\mathbb{K}} \mathbf{P}(\mathbf{V})\) 或 \(\dim \mathbf{P}(\mathbf{V})\) 。因此，维数为 -1 的射影空间是空集，维数为 0 的射影空间是单个点。维数为 1（相应地，2）的射影空间称为射影直线（相应地，射影平面）。

此后我们仅考虑左射影空间。

\subsection{齐次坐标}

设 V 是 K 上一个维数为 \(n + 1\) 的有限维向量空间， \(\mathbf{P}(\mathbf{V})\) 是由 V 导出的 n 维射影空间， \((e_{i})_{0 \leqslant i \leqslant n}\) 是 V 的一个基。设 \(\pi\) 表示 V - \{0\} 到商集 \(\mathbf{P}(\mathbf{V})\) 上的典范映射。对每一点 \(x = \sum_{i=0}^{n} \xi_i e_i\) （属于 \(V - \{0\}\) ）， \((\xi_0, \xi_1, \ldots, \xi_n)\) 称为点 \(\pi(x)\) 关于基 \((e_i)\) （属于 \(V\) ）的齐次坐标系。因此，每个系 \((\xi_i)\) （由 \(n + 1\) 个 \(K\) 中不全为零的元素组成）都是 \(P(V)\) 的一个点关于 \((e_i)\) 的一个齐次坐标系；要使两个这样的系 \((\xi_i), (\xi_i')\) 是 \(P(V)\) 的同一点关于同一基 \((e_i)\) 的齐次坐标系，必要且充分条件是：存在元素 \(\lambda \neq 0\) （在 \(K\) 中），使得 \(\xi_i' = \lambda \xi_i\) 对 \(0 \leqslant i \leqslant n\) 成立。

此定义可立即推广到 V 为无限维的情形。

给定 V 的另一个基 \((\bar{e}_i)\) ，使得 \(e_i = \sum_{j=0}^{n} \alpha_{ij} \bar{e}_j\) ( \(0 \leqslant i \leqslant n\) )，以及一个齐次坐标系 \((\xi_i)\) ，即 \(\pi(x)\) 关于基 \((e_i)\) 的坐标系；要使系 \((\bar{\xi}_i)\) （由 \(n+1\) 个 K 的元素组成）是 \(\pi(x)\) 关于基 \((\bar{e}_i)\) 的一个齐次坐标系，必要且充分条件是：存在 \(\lambda \neq 0\) （在 K 中），使得

\[
\lambda \bar {\xi} _ {i} = \sum_ {j = 0} ^ {n} \xi_ {j} \alpha_ {j i} \quad \text { for } 0 \leqslant i \leqslant n.
\]

特别地，若 \(e_i = \gamma_i \bar{e}_i\) ，其中 \(\gamma_i \neq 0\) ( \(0 \leqslant i \leqslant n\) )，则 \(\bar{\xi}_i = \mu \xi_i \gamma_i\) ，其中 \(\mu \neq 0\) 。

\subsection{射影线性簇}

设 W 是向量空间 V 的一个向量子空间； \(W - \{0\}\) 在由 V 导出的射影空间 \(\mathbf{P}(V)\) 中的典范像称为射影线性簇（若无混淆之虞，则简称为线性簇）；由于等价关系 \(\Delta(W)\) （定义在 \(W - \{0\}\) 上）是由关系 \(\Delta(V)\) 诱导的， \(W - \{0\}\) 在 \(\mathbf{P}(V)\) 中的像这一射影线性簇可以与由 W 导出的射影空间 \(\mathbf{P}(W)\) 等同起来，因而我们可以谈论这种簇的维数。在射影空间 \(\mathbf{P}(V)\) 中，V 的（去掉原点的）超平面的典范像是一个线性簇，称为射影超平面（或简称超平面）；若 \(\mathbf{P}(V)\) 是 n 维有限维的，则 \(\mathbf{P}(V)\) 中的超平面是维数为 n - 1 的线性簇。

关于向量空间的向量子空间的每一个命题，都可以转化为关于射影线性簇的命题。例如，如果一个射影空间 \(\mathbf{P}(\mathbf{V})\) 是有限维的，维数为 n，且 \((e_{i})_{0\leqslant i\leqslant n}\) 是 V 的一组基，那么每一个维数为 r 的线性簇 \(\mathbf{L}\subset\mathbf{P}(\mathbf{V})\) 都可以由 n-r 个齐次线性方程组成的方程组来定义

\[
\sum_ {i = 0} ^ {n} \xi_ {i} \alpha_ {i j} = 0 \quad (1 \leqslant j \leqslant n - r)\tag{4}
\]

在齐次坐标 \(\xi_{i}\) ( \(0 \leqslant i \leqslant n\) )（即 \(\mathbf{P}(\mathbf{V})\) 的点关于基 \((e_{i})\) 的齐次坐标）之间，其中 (4) 的左边是 V 上线性无关的形式。特别地，一个射影超平面由系数不全为零的单个齐次线性方程定义。反之， \(\mathbf{P}(\mathbf{V})\) 中满足关于 \(\xi_{i}\) 的任意齐次线性方程组的点构成一个线性簇 L；若所考虑的方程组由 \(k \leqslant n + 1\) 个方程组成，则 L 的维数 \(\geqslant n - k\) 。

 \(\mathbf{P}(\mathbf{V})\) 中线性簇的任何交都是一个线性簇；对 \(\mathbf{P}(\mathbf{V})\) 的每个子集 A，存在包含 A 的最小线性簇 L；它称为由 A 生成的线性簇，而 A 称为 L 的一个生成系。

如果 W 是 V 的由 \(\pi^{1}(A)\) ，则 L = P(W)。

若 \(\mathbf{L}\) 和 \(\mathbf{M}\) 是 \(\mathbf{P}(\mathbf{V})\) 中任意两个线性簇， \(\mathbf{N}\) 是由 \(\mathbf{L} \cup \mathbf{M}\) 生成的簇，则（§7，第 3 号，命题 4 的推论 3）

\[
\dim \mathbf {L} + \dim \mathbf {M} = \dim (\mathbf {L} \cap \mathbf {M}) + \dim \mathbf {N}.\tag{5}
\]

特别是，如果 \(\mathbf{P}(\mathbf{V})\) 是有限维的，且 \(\dim \mathbf{L} + \dim \mathbf{M} \geqslant \dim \mathbf{P}(\mathbf{V})\) ，由（5）可知 \(\mathbf{L} \cap \mathbf{M}\) 是非空的。

设 \((x_{i}), (y_{i})\) 是向量空间 V 中具有同一指标集的两族点，满足 \(y_{i} = \lambda_{i}x_{i}\) ，其中对所有 i 有 \(\lambda_{i}^{*} \neq 0\) 。若族 \((x_{i})\) 是自由的，则 \((y_{i})\) 也是自由的，反之亦然；这时称 P(V) 的点族 \(\pi(x_{i})\) 是射影自由的（或简称自由的）。这等价于说：对每个指标 K，点 \(\pi(x_{K})\) 不属于由诸 \(\pi(x_{i})\) （ \(i \neq K\) ）生成的线性簇。P(V) 中不是射影自由的点族称为射影相关的（或简称相关的）。

要使点的族 \((x_{i})\) （属于 \(\mathbf{V} - \{0\}\) ）满足：族 \((\pi (x_{i}))\) 射影自由且生成 \(\mathbf{P}(\mathbf{V})\) ，必要且充分条件是 \((x_{i})\) 是 \(\mathbf{V}\) 的一个基。若 \(\mathbf{P}(\mathbf{V})\) 的维数为 \(n\) ，则这样一个族中元素的个数是 \(n + 1\) 。注意，给定这样一个族 \((\pi (x_{i}))\) （在 \(\mathbf{P}(\mathbf{V})\) 中）并不确定 \(\mathbf{P}(\mathbf{V})\) 的一个给定点关于基 \((y_{i})\) （属于 \(\mathbf{V}\) ）的齐次坐标（即使允许相差一个左因子），其中 \(\pi (y_{i}) = \pi (x_{i})\) 对所有 \(\iota\) 成立（参见第 6 号）。

\subsection{仿射空间的射影完备化}

设 V 是域 K 上的一个（左）向量空间，考虑 K 上的向量空间 \(K_{s} \times V\) ；射影空间 \(\mathbf{P}(\mathbf{K}_{s} \times \mathbf{V})\) 称为与向量空间 V 典范相伴的射影空间。若 V 的维数为 n，则 \(\mathbf{P}(\mathbf{K}_{s} \times \mathbf{V})\) 有相同的维数 n。考虑 \(K_{s} \times V\) 中的仿射超平面 \(V_{1} = \{1\} \times V\) ，它的方向（第 3 号）是子空间 \(V_{0} = \{0\} \times V\) ；若 \(K_{s} \times V\) 的一条（过 0 的）直线不含于 \(V_{0}\) ，则它包含一个点 \((\alpha, x)\) ，其中 \(\alpha \neq 0\) 且 \(x \in V\) ，因而它也包含点 \(\alpha^{-1}(\alpha, x) = (1, \alpha^{-1}x)\) （属于 \(V_{1}\) ）；其逆是直接的，于是可以看出，在 \(V_{1}\) 的点与 \(\mathbf{K}_s \times \mathbf{V}\) 中不含于 \(\mathbf{V}_0\) 的（过 0 的）直线之间存在一一对应，后者的每一条与 \(\mathbf{V}_1\) 交于且仅交于一点。由此可知，映射 \(x \mapsto \phi(x) = \pi(1, x)\) 是 \(\mathbf{V}\) 到射影空间 \(\mathbf{P}(\mathbf{K}_s \times \mathbf{V})\) 中的一个单射（称为典范单射）； \(\mathbf{V}\) 常与此单射下的像等同。 \(\phi(\mathbf{V})\) 在 \(\mathbf{P}(\mathbf{K}_s \times \mathbf{V})\) 中的补集是射影超平面 \(\mathbf{P}(\mathbf{V}_0)\) ，称为 \(\mathbf{P}(\mathbf{K}_s \times \mathbf{V})\) （或按滥用的说法， \(\mathbf{V}\) 的）的无穷远超平面；它的点也称为 \(\mathbf{P}(\mathbf{K}_s \times \mathbf{V})\) （或 \(\mathbf{V}\) ）的无穷远点。若 \((a_t)\) 是 \(\mathbf{V}\) 的一个基，且在 \(\mathbf{K}_s \times \mathbf{V}\) 中取由元素 \(e_t = (0, a_t)\) 与元素 \(e_\omega = (1, 0)\) 组成的基，则 \(\mathbf{P}(\mathbf{K}_s \times \mathbf{V})\) 中的无穷远点是那些指标 \(\omega\) 的齐次坐标为 0 的点。

设 M 是 V 中的一个仿射线性簇（第 3 号），D 是它的方向；M 的典范像 \(\phi(M)\) （在 \(\mathbf{P}(\mathbf{K}_s \times \mathbf{V})\) 中）含于典范像 \(\overline{\mathbf{M}} = \pi(\mathbf{M}_2)\) ，它是向量子空间 \(\mathbf{M}_2\) （ \(\mathbf{K}_s \times \mathbf{V}\) 中由仿射簇 \(\mathbf{M}_1 = \{1\} \times \mathbf{M}\) （ \(\mathbf{K}_s \times \mathbf{V}\) 中的）所生成的）的典范像。更精确地说，若 \((a_t)\) 是 M 的一个生成 M 的仿射自由系，则诸元素 \((1, a_t)\) 构成 \(\mathbf{M}_2\) 的一个基，从而 \(\overline{\mathbf{M}}\) 正是由 \(\phi(M)\) 生成的射影线性簇；若 M 是有限维的，则 \(\overline{\mathbf{M}}\) 与 M 有相同的维数。 \(\phi(M)\) 在 \(\overline{\mathbf{M}}\) 中的补集是 \(\overline{\mathbf{M}}\) 与无穷远超平面的交，且等于典范像 \(\pi(\mathbf{M}_0)\) ，其中 \(\mathbf{M}_0 = \{0\} \times \mathbf{D}\) 。

反之，设 N 是一个不含于无穷远超平面的射影线性簇，并设 \(\mathbf{R} = \overline{\pi}^{-1}(\mathbf{N})\) ； \(R \cap V_{1}\) 是 \(K \times V\) 的一个仿射线性簇，形如 \(\{1\} \times M\) ，其中 M 是 V 的一个仿射线性簇，并且立即可看出，N 是仿射线性簇 \(\overline{M}\) ，它由 \(\phi(M)\) 生成。

因此，在 V 的仿射线性簇与 \(\mathbf{P}(\mathbf{K}_{s} \times \mathbf{V})\) 中不含于无穷远超平面的射影线性簇之间存在一一对应；V 的两个仿射线性簇平行，必要且充分条件是它们生成的射影线性簇与无穷远超平面有相同的交（这有时表述为：所考虑的两个仿射线性簇具有相同的无穷远点）。

\subsection{有理函数的延拓}

若把第 8 号的结果应用于维数为 1 的向量空间 \(\mathbf{V} = \mathbf{K}_s\) ，则可看出存在一个典范单射 \(\phi\) ，把 \(\mathbf{K}_s\) 映入射影直线 \(\mathbf{P}_1(\mathbf{K}) = \mathbf{P}(\mathbf{K}_s \times \mathbf{K}_s)\) ；对所有 \(\xi \in \mathbf{K}\) ， \(\phi(\xi)\) 是这样的点：其齐次坐标 \((1, \xi)\) 是关于典范基（ \(\S 1\) ，第 11 号）——即 \(\mathbf{K}_s \times \mathbf{K}_s\) 的典范基——取的。 \(\phi(\mathbf{K})\) 在 \(\mathbf{P}_1(\mathbf{K})\) 中的补集由单独一个点组成，它关于上述基的齐次坐标为 \((0, 1)\) ；它称为“无穷远点”。 \(\mathbf{P}_1(\mathbf{K})\) 也称为与 \(\mathbf{K}\) 相伴的射影域，记作 \(\tilde{\mathbf{K}}\) ， \(\tilde{\mathbf{K}}\) 中的无穷远点记作 \(\infty\) 。

特别考虑 K 为交换域的情形，并设 \(f \in \mathbf{K}(\mathbf{X})\) 是 \(\mathbf{K}\) 上一个未定元的有理函数（IV，§4）；若 \(f \neq 0\) ，则存在唯一的表达式 \(f = \alpha p/q\) ，其中 \(\alpha \in \mathbf{K}^*\) ， \(p\) 与 \(q\) 是两个互素的首一多项式（VII，§1）；设 \(m\) 与 \(n\) 分别是它们的次数，并设 \(r = \sup(m, n)\) 。我们记

\[
p _ {1} (\mathrm{T}, \mathrm{X}) = \mathrm{T} ^ {r} p (\mathrm{X} / \mathrm{T}), \quad q _ {1} (\mathrm{T}, \mathrm{X}) = \mathrm{T} ^ {r} q (\mathrm{X} / \mathrm{T});
\]

\(p_{1}\) 与 \(q_{1}\) 是 K 上两个次数为 r 的齐次多项式，满足 \(p(\mathbf{X}) = p_{1}(1, \mathbf{X})\) , \(q(\mathbf{X}) = q_{1}(1, \mathbf{X})\) 。因此，对每个元素 \(\xi \in K\) （它不是 \(q(\mathbf{X})\) 的零点）， \(f(\xi) = \alpha p(\xi)/q(\xi)\) 有定义，且我们可以写出

\[
f (\xi) = \alpha p _ {1} (1, \xi) / q _ {1} (1, \xi) = \alpha p _ {1} (\lambda , \lambda \xi) / q _ {1} (\lambda , \lambda \xi)
\]

对所有 \(\lambda \neq 0\) （在 K 中）成立。于是考虑映射

\[
(\eta , \xi) \mapsto (q _ {1} (\eta , \xi), p _ {1} (\eta , \xi))
\]

（ \(K^{2}\) 到自身的映射）；它与等价关系 \(\Delta(\mathbf{K}^{2})\) 相容，因而在过渡到商时定义了一个映射 \(\tilde{f}\) （ \(\tilde{K}\) 到自身），它在有理函数有定义的点处与 \(\xi\mapsto f(\xi)\) 一致；按语言滥用的说法，称 \(\tilde{f}\) 为 f 到 \(\tilde{K}\) 的典范延拓。

例如，若 \(f = 1 / \mathbf{X}\) ，则 \(\tilde{f}(0) = \infty\) 且 \(\tilde{f}(\infty) = 0\) ；若

\[
f = (a \mathbf {X} + b) / (c \mathbf {X} + d)
\]

其中 \(ad - bc \neq 0\) ，则 \(\tilde{f}(-d/c) = \infty\) ， \(\tilde{f}(\infty) = a/c\) （若 \(c \neq 0\) ）， \(\tilde{f}(\infty) = \infty\) （若 \(c = 0\) ）。若 \(f = a_0\mathbf{X}^n + \cdots + a_n\) 是一个次数为 \(n > 0\) 的多项式，则 \(\tilde{f}(\infty) = \infty\) 。

\subsection{射影线性映射}

设 V、 \(V'\) 是域 K 上的两个左向量空间，f 是 V 到 \(V'\) 的一个线性映射， \(N = \bar{f}^{-1}(0)\) 是它的核。立即可见，V 中一条不含于 N 的（过 0 的）直线在 f 下的像是 \(V'\) 中一条（过 0 的）直线；因此，在过渡到商时，f 定义了一个映射 g （从 \(\mathbf{P}(\mathbf{V}) - \mathbf{P}(\mathbf{N})\) 到 \(\mathbf{P}(\mathbf{V}')\) ）。这样的映射称为射影线性映射（或简称射影映射）；尽管它定义在 \(\mathbf{P}(\mathbf{V}) - \mathbf{P}(\mathbf{N})\) 上而不是 \(\mathbf{P}(\mathbf{V})\) 上（当 \(N \neq \{0\}\) 时），我们仍按语言滥用的说法称 g 是 \(\mathbf{P}(\mathbf{V})\) 到 \(\mathbf{P}(\mathbf{V}')\) 的一个射影映射。射影线性簇 \(\mathbf{P}(\mathbf{N})\) （g 在其上无定义）称为 g 的中心。

注意，当 \(g\) 定义在整个 \(\mathbf{P}(\mathbf{V})\) 上时（即当 \(\mathbf{N} = \{0\}\) 时）， \(g\) 是 \(\mathbf{P}(\mathbf{V})\) 到 \(\mathbf{P}(\mathbf{V}')\) 的一个单射。

当基 \((a_{\lambda})_{\lambda\in\mathbf{L}}, (b_{\mu})_{\mu\in\mathbf{M}}\) 分别在 V 与 \(V'\) 中给定时， \(\mathbf{P}(\mathbf{V})\) 到 \(\mathbf{P}(\mathbf{V}')\) 的一个射影映射把 \(\mathbf{P}(\mathbf{V})\) 中具有齐次坐标 \(\xi_{\lambda}\) ( \(\lambda \in \mathbf{L}\) ) 的点映到 \(\mathbf{P}(\mathbf{V}')\) 中的一点，该点具有形如以下的齐次坐标系 \(\eta_{\mu}\) ( \(\mu \in \mathbf{M}\) )

\[
\eta_ {\mu} = \sum_ {\lambda \in L} \xi_ {\lambda} \alpha_ {\lambda \mu} \quad (\alpha_ {\lambda \mu} \in K).\tag{6}
\]

 \(g\) 的中心是由以下方程定义的线性簇

\[
\sum_ {\lambda \in L} \xi_ {\lambda} \alpha_ {\lambda \mu} = 0 \quad (\mu \in M).
\]

如果 C 是 g 的中心，且 M 是 \(\mathbf{P}(\mathbf{V})\) 的一个线性簇，则 \(\mathbf{M} - (\mathbf{M} \cap \mathbf{C})\) 在 g 下的像是 \(\mathbf{P}(\mathbf{V}')\) 的一个线性簇，（按语言滥用）记作 \(g(\mathbf{M})\) 。则

\[
\dim g (\mathbf {M}) + \dim (\mathbf {M} \cap \mathbf {C}) + 1 = \dim \mathbf {M}\tag{7}
\]

（§7，第 4 号，公式 (12)）。若 \(\mathbf{M}'\) 是 \(\mathbf{P}(\mathbf{V}')\) 的一个线性簇，则 \(\bar{g}^{-1}(\mathbf{M}') \cup \mathbf{C}\) 是 \(\mathbf{P}(\mathbf{V})\) 的一个线性簇，并且

\[
\dim (\bar {g} ^ {- 1} (\mathbf {M} ^ {\prime}) \cup \mathbf {C}) = \dim \mathbf {C} + \dim (\mathbf {M} ^ {\prime} \cap g (\mathbf {P} (\mathbf {V}))) + 1.\tag{8}
\]

我们按语言滥用的说法，称 \(\bar{g}^{-1}(\mathbf{M}') \cup \mathbf{C}\) 为 \(\mathbf{M}'\) 在 \(g\) 下的逆像。

由于线性映射在基 \((e_{t})\) （属于 V ）上取的值可以在 \(V'\) 中任意选取，可见存在 \(\mathbf{P}(\mathbf{V})\) 到 \(\mathbf{P}(\mathbf{V}')\) 的一个射影线性映射，它在点 \(\pi(e_{t})\) 处取任意的值。但是（即使当 g 处处有定义时）给出 \(g(\pi(e_{t}))\) 也不能唯一地确定 g（习题 10）。

两个双射的射影映射的复合仍是一个射影映射；这样的双射的逆映射也是射影映射。射影空间 \(\mathbf{P}(\mathbf{V})\) 到自身上的双射射影映射因此构成一个群，称为 \(\mathbf{P}(\mathbf{V})\) 的射影群，记作 \(\mathbf{PGL}(\mathbf{V})\) ；我们用 \(\mathbf{PGL}_{n}(\mathbf{K})\) 或 \(\mathbf{PGL}(n,\mathbf{K})\) 代替 \(\mathbf{PGL}(\mathbf{K}_{s}^{n})\) 。

注记。在射影空间 \(\mathbf{P}(\mathbf{V})\) （在域 \(\mathbf{K}\) 上）中，设 \(\mathbf{H} = \mathbf{P}(\mathbf{W})\) 是一个超平面。存在一个双射线性映射 \(f\) （从 \(\mathbf{V}\) 映到 \(\mathbf{K}_s \times \mathbf{W}\) 上），使得 \(f(\mathbf{W}) = \mathbf{W}\) ；设 \(g\) 是由 \(f\) 通过过渡到商而得到的射影映射。已经看到（第 8 号）， \(\mathbf{P}(\mathbf{W})\) 在 \(\mathbf{P}(\mathbf{K}_s \times \mathbf{W})\) 中的补集可以等同于一个仿射空间，其平移空间是 \(\mathbf{W}\) 。当 \(\mathbf{P}(\mathbf{V})\) 与 \(\mathbf{P}(\mathbf{K}_s \times \mathbf{W})\) 借助 \(g\) 等同时，就说取 \(\mathbf{H}\) 为 \(\mathbf{P}(\mathbf{V})\) 中的无穷远超平面；于是 \(\mathbf{H}\) 在 \(\mathbf{P}(\mathbf{V})\) 中的补集等同于一个仿射空间，其平移空间是 \(\mathbf{W}\) 。

\subsection{射影空间结构}

给定一个集合 E 和一个域 K，E 上关于域 K 的（左）射影空间结构由给出一个非空集合 \(\Phi\) ——它由射影空间 \(\mathbf{P}(\mathbf{K}_{s}^{\mathrm{(E)}})\) 的一些子集到 E 上的双射组成——来定义，且满足以下公理：

\((\mathbf{EP}_{\mathrm{I}})\) 每个映射 \(f \in \Phi\) 的定义集是 \(\mathbf{P}(\mathbf{K}_s^{(\mathbf{E})})\) 的一个线性簇。

\((\mathrm{EP}_{\Pi})\) 对每一对元素 \(f, g\) （属于 \(\Phi\) ，分别定义在线性簇 \(\mathbf{P}(\mathbf{V})\) 与 \(\mathbf{P}(\mathbf{W})\) 上），双射 \(h = \overline{g}^{-1} \circ f\) （ \(\mathbf{P}(\mathbf{V})\) 到 \(\mathbf{P}(\mathbf{W})\) 上）是一个射影映射。

\((\mathbf{EP}_{\mathrm{III}})\) 反之，若 \(f \in \Phi\) 定义在线性簇 \(\mathbf{P}(\mathbf{V})\) 上，而 \(h\) 是 \(\mathbf{P}(\mathbf{V})\) 到一个线性簇 \(\mathbf{P}(\mathbf{W}) \subset \mathbf{P}(\mathbf{K}_s^{\mathbf{(E)}})\) 上的双射射影映射，则 \(f \circ h^{-1} \in \Phi\) 。

设 E 是一个集合， \((\mathbf{V}_{\lambda})_{\lambda \in L}\) 是 K 上的一族向量空间，并且假设对每个 \(\lambda \in L\) 给定了一个双射 \(f_{\lambda}\) （从 \(\mathbf{P}(\mathbf{V}_{\lambda})\) 到 E 上），使得对每一对有序指标 \(\lambda, \mu, f_{\lambda} \circ f_{\mu}\) 而言，它是 \(\mathbf{P}(\mathbf{V}_{\mu})\) 到 \(\mathbf{P}(\mathbf{V}_{\lambda})\) 的一个射影映射。于是我们可以如下在 E 上定义关于 K 的射影空间结构：设 \((e_{t})_{t \in I}\) 是空间 \(V_{\lambda}\) 的一个基，记 \(a_{t} = f_{\lambda}(\pi(e_{t}))\) ；设 \(b_{t}\) 是指标为 \(a_{t}\) 的元素（在 \(\mathbf{K}_{s}^{(\mathbf{E})}\) 的典范基中）（§1，第 11 号）。关系 \(t \neq k\) 蕴含 \(b_{t} \neq b_{k}\) ，这是因为 \(f_{\lambda}\) 是双射这一假设；因此诸 \(b_{t}\) 构成向量子空间 \(W_{0}\) （ \(K_{s}^{(\mathbf{E})}\) 的）的一个基，从而存在 \(\mathbf{P}(\mathbf{W}_{0})\) 到 \(\mathbf{P}(\mathbf{V}_{\lambda})\) 上的一个双射射影映射 h ，使得 \(h(\pi(b_{t})) = \pi(e_{t})\) 对所有 \(t \in I\) 成立。若取 \(\Phi\) 为所有双射射影映射 \(f_{\lambda} \circ h \circ g^{-1}\) 的集合，其中 g 遍历所有双射射影映射 \(\mathbf{P}(\mathbf{W}) \subset \mathbf{P}(\mathbf{K}_{s}^{(\mathbf{E})})\) 的集合，则立即可验证 \(\Phi\) 满足公理 \((\mathrm{EP}_{\mathrm{I}}), (\mathrm{EP}_{\mathrm{II}})\) 与 \((\mathrm{EP}_{\mathrm{III}})\) 。此外显然 \(\Phi\) 既不依赖于指标 \(\lambda \in L\) 的选取，也不依赖于基 \((e_{t})\) （在 \(V_{\lambda}\) 中）的选取，也不依赖于 h 的选取。

特别地（取 L 由单个元素组成），由向量空间 V 导出的每个射影空间 \(\mathbf{P}(\mathbf{V})\) 因此具有本号所给定义意义下的一个完全确定的“射影空间结构”。因此，任何具有射影空间结构的集合都可以称为射影空间。

\[
f \in \Phi
\]

\[
\mathbf {P} (\mathbf {V}) \subset \mathbf {P} \left(\mathrm{K} _ {\mathrm{s}} ^ {(\mathrm{E})}\right)
\]

\(f(\mathbf{M})\) 是 \(\mathbf{P}(\mathbf{V})\) 中第 7 号意义下的线性簇（此性质于是对所有 \(f \in \Phi\) 成立）。由此可知，射影空间中的每个线性簇都典范地具有一个射影空间结构。

称射影空间 \(\mathbf{E}\) 的维数为 \(n\) ，如果对所有 \(f\in\Phi,f(\mathbf{E})\) 都是维数为 \(n\) 的线性簇（只需对一个映射 \(f\in\Phi\) 成立即可）。

\section{矩阵}

\subsection{矩阵的定义}

定义 1. 设 I, K, H 为三个集合；一个 (I, K) 型、元素属于 H 的矩阵（或 H 上的 (I, K) 型矩阵）是指 H 的元素的任意族 \(M = (m_{\iota\kappa})_{(\iota, \kappa) \in I \times K}\) ，其指标集为乘积 \(I \times K\) 。对所有 \(\iota \in I\) ，族 \((m_{\iota\kappa})_{\kappa \in K}\) 称为 \(M\) 的指标为 \(\iota\) 的行；对所有 \(\kappa \in K\) ，族 \((m_{\iota\kappa})_{\iota \in I}\) 称为 \(M\) 的指标为 \(\kappa\) 的列。

如果 I（相应地 K）是有限的，则称 M 为行数（相应地列数）有限的矩阵。H 上 (I, K) 型矩阵的集合与乘积 \(H^{I \times K}\) 等同。

名称“行”和“列”源于以下事实：在 I 和 K 是区间 \([1,p]\) , \([1,q]\) 的 N 的情况下，矩阵的元素被视为排列在一个矩形阵列中，该阵列有 p 行（水平排列）和 q 列（垂直排列）：

\[
\left( \begin{array}{c c c c c} m _ {1 1} & m _ {1 2} & \ldots & m _ {1 q} \\ m _ {2 1} & m _ {2 2} & \ldots & m _ {2 q} \\ \cdot & \cdot & \cdot & \cdot & \cdot \\ m _ {p 1} & m _ {p 2} & \ldots & m _ {p q} \end{array} \right)
\]

当p和q是足够小的显式整数，使得这样做切实可行时，约定上述阵列是有效表示所讨论矩阵的符号；这种记号使我们无需使用指标，因为一个元素的指标由它在阵列中的位置决定；例如，当我们谈到矩阵

\[
\left( \begin{array}{c c c} a & b & c \\ d & e & f \end{array} \right)
\]

我们指的是矩阵 \((m_{ij})_{1\leqslant i\leqslant 2,1\leqslant j\leqslant 3}\) ，使得

\[
m _ {1 1} = a, m _ {1 2} = b, m _ {1 3} = c, m _ {2 1} = d, m _ {2 2} = e, m _ {2 3} = f.
\]

代替类型为 \(([1, p], [1, q])\) 的矩阵的说法，我们也说 \((p, q)\) 型矩阵，或具有 \(p\) 行 \(q\) 列的矩阵（若不致引起混淆）；H 上 \((p, q)\) 型矩阵的集合有时记作 \(\mathbf{M}_{p,q}(\mathrm{H})\) 。

H 上的每个矩阵，若其指标集 I、K 中有一个为空，则与 H 的元素的空族相同；它也称为空矩阵。当 \(I = \{i_{0}\}\) （相应地 \(K = \{k_{0}\}\) ）是仅含一个元素的集合时，M 称为行矩阵（相应地列矩阵），此时行（相应地列）指标可在记号中省去；当 I 与 K 都是单元素集时，(I, K) 型矩阵常被等同于这个矩阵中的唯一元素。

\(M' = (m_{\iota \kappa})_{(\iota, \kappa) \in J \times L}\) （矩阵 \(M = (m_{\iota \kappa})_{(\iota, \kappa) \in I \times K}\) 的一个子族），若其指标集是子集 \(J\) （ \(I\) 的）与子集 \(L\) （ \(K\) 的）的乘积，称为矩阵 \(M\) 的子矩阵；说它是从 \(M\) 中删去指标为 \(\iota \notin J\) 的行与指标为 \(\kappa \notin L\) 的列而得到的；反之，说 \(M\) 是由 \(M'\) 通过加边——加上指标为 \(\iota \notin J\) 的行与指标为 \(\kappa \notin L\) 的列——而得到的。

定义 2. 矩阵 \(M = (m_{\iota \kappa})_{(\iota, \kappa) \in \mathbf{I} \times \mathbf{K}}\) 的转置，记作 \({}^tM\) ，是矩阵 \((m_{\kappa \iota}')_{(\kappa, \iota) \in \mathbf{K} \times \mathbf{I}}\) （ \(\mathbf{H}\) 上的），它由 \(m_{\kappa \iota}' = m_{\iota \kappa}\) （对所有 \((\iota, \kappa) \in \mathbf{I} \times \mathbf{K}\) ）给出。

由该定义可知，(I, K)型矩阵的转置是(K, I)型矩阵，并且

\[
{ } ^ { t } ( { } ^ { t } M ) = M .\tag{1}
\]

\subsection{交换群上的矩阵}

设 G 是一个交换群（写成加法）。带给定指标集 I, K 的 G 上矩阵的集合具有交换群结构，因为它是 \(I \times K\) 到 G 的映射集合；这个群写成加法，于是若 \(M = (m_{\iota\kappa})\) 与 \(M' = (m_{\iota\kappa}')\) 是它的两个元素，则

\[
M + M ^ {\prime} = (m _ {\iota \kappa} + m _ {\iota \kappa} ^ {\prime});
\]

因此，这个群的单位元是所有元素均为零的矩阵（称为零矩阵）。显然，

\[
{ } ^ { t } ( M + M ^ { \prime } ) = { } ^ { t } M + { } ^ { t } M ^ { \prime } .\tag{2}
\]

两个矩阵的和仅在这两个矩阵的行索引集和列索引集相同时才有定义。

设 \(\mathbf{H}'\) , \(\mathbf{H}''\) 是两个集合， \(\mathbf{G}\) 是一个交换群（写成加法），而 \(f: (h', h'') \mapsto h'h''\) 是 \(\mathbf{H}' \times \mathbf{H}''\) 到 \(\mathbf{G}\) 的一个映射。给定两个矩阵

\[
M ^ {\prime} = \left(m _ {i k} ^ {\prime}\right) _ {(i, k) \in I \times K}, \quad M ^ {\prime \prime} = \left(m _ {k l} ^ {\prime \prime}\right) _ {(k, l) \in K \times L}
\]

分别位于 H' 与 H'' 上，使得 M' 的列的指标集 K 是有限的且等于 M'' 的行的指标集，则 M' 与 M'' 经由 f 的乘积，记作 M'M'' 或 f(M', M'')，是矩阵

\[
\left(\sum_ {k \in \mathbf {K}} m _ {i k} ^ {\prime} m _ {k l} ^ {\prime \prime}\right) _ {(i, l) \in \mathbf {I} \times \mathbf {L}}\tag{3}
\]

（ G 上的）。

上述定义假定 \(M'\) 的列的指标集等于 \(M''\) 的行的指标集；特别地，乘积 \(M''M'\) 没有意义，若 \(I \neq L\) 。在公式 (3) 中，同一

 \(M'\) 的同一行的元素以 \(M''\) 的同一列的元素右乘而出现；称这种乘法是按“行乘列”进行的。

设 \(f^0\) 是映射 \((h'', h') \mapsto h'h''\) （从 \(H'' \times H'\) 到 G）；由定义立即可得

\[
{ } ^ { t } ( M ^ { \prime } M ^ { \prime \prime } ) = { } ^ { t } M ^ { \prime \prime } . { } ^ { t } M ^ { \prime }\tag{4}
\]

其中左边（相应地右边）的乘积是经由 \(f\) （相应地经由 \(f^0\) ）计算的。

当 \(\mathbf{H}'\) 与 \(\mathbf{H}''\) 本身是交换群（写成加法）且 \(f\) 是 \(\mathbf{Z}\) -双线性（ \(\S 3\) ，第 1 号）时，分配律公式

\[
\left\{ \begin{array}{l} (M ^ {\prime} + N ^ {\prime}) M ^ {\prime \prime} = M ^ {\prime} M ^ {\prime \prime} + N ^ {\prime} M ^ {\prime \prime} \\ M ^ {\prime} (M ^ {\prime \prime} + N ^ {\prime \prime}) = M ^ {\prime} M ^ {\prime \prime} + M ^ {\prime} N ^ {\prime \prime} \end{array} \right.\tag{5}
\]

立即得到验证，指标集使得其中出现的和与积均有定义。

现在设 \(H_{1}\) , \(H_{2}\) , \(H_{3}\) , \(H_{12}\) , \(H_{23}\) 与 H 是交换群（写成加法）， \(f_{12}:H_{1} \times H_{2} \to H_{12}\) , \(f_{23}:H_{2} \times H_{3} \to H_{23}\) 是映射，而

\[
f _ {3}: \mathrm{H} _ {1 2} \times \mathrm{H} _ {3} \rightarrow \mathrm{H}, \quad f _ {1}: \mathrm{H} _ {1} \times \mathrm{H} _ {2 3} \rightarrow \mathrm{H}
\]

Z-双线性映射；进一步假设，对所有 \(x_{i} \in H_{i}\) （i = 1, 2, 3）

\[
f _ {3} (f _ {1 2} (x _ {1}, x _ {2}), x _ {3}) = f _ {1} (x _ {1}, f _ {2 3} (x _ {2}, x _ {3}))
\]

（它也可像上面那样写成 \((x_{1}x_{2})x_{3} = x_{1}(x_{2}x_{3}))\) ；于是，若 \(M' = (m'_{rs})\) , \(M'' = (m''_{st})\) , \(M''' = (m''_{tu})\) 分别是 \(\mathrm{H}_1\) , \(\mathrm{H}_2\) , \(\mathrm{H}_3\) 上的矩阵，

\[
(M ^ {\prime} M ^ {\prime \prime}) M ^ {\prime \prime \prime} = M ^ {\prime} (M ^ {\prime \prime} M ^ {\prime \prime \prime})\tag{6}
\]

（当两边的乘积（分别经由 \(f_{12}\) , \(f_{3}\) , \(f_{23}\) 与 \(f_{1}\) 计算）都有定义时）成立；因为

\[
\begin{array}{r l} \sum_ {t} \left(\sum_ {s} m _ {r s} ^ {\prime} m _ {s t} ^ {\prime \prime}\right) m _ {t u} ^ {\prime \prime \prime} & = \sum_ {t} \sum_ {s} (m _ {r s} ^ {\prime} m _ {s t} ^ {\prime \prime}) m _ {t u} ^ {\prime \prime \prime} = \sum_ {s} \sum_ {t} m _ {r s} ^ {\prime} (m _ {s t} ^ {\prime \prime} m _ {t u} ^ {\prime \prime \prime}) \\ & = \sum_ {s} m _ {r s} ^ {\prime} \left(\sum_ {t} m _ {s t} ^ {\prime \prime} m _ {t u} ^ {\prime \prime \prime}\right) \end{array}
\]

根据所作的假设，

（6）式的两边也记作 \(M'M''M'''\) 。对于多于三个因子的乘积，也采用类似的约定。

注记。上述公式可推广到更一般的情形。确切地说：

\begin{enumerate}
\def\labelenumi{(\alph{enumi})}
\item
假设 \(H = \bigcup_{(\iota, \kappa) \in I \times K} G_{\iota\kappa}\) ，其中每个 \(G_{\iota\kappa}\) 是写成加法的交换群；那么和 \(M + M'\) 可以在如下情形定义：对每个有序对 \((\iota, \kappa)\) ， \(m_{\iota\kappa} \in G_{\iota\kappa}\) 且 \(m'_{\iota\kappa} \in G_{\iota\kappa}\) 。
\item
  设 I、K、L 为三个集合，其中 K 假定为有限集，且设 \(H' = \bigcup_{(i, k) \in I \times K} H'_{ik}\) , \(H'' = \bigcup_{(k, l) \in K \times L} H''_{kl}\) , \(H = \bigcup_{(i, l) \in I \times L} H_{il}\) 是三个集合；假设每一个 \(H_{il}\) 是一个交换群，以加法形式书写，并且对于每个三元组 \((i, k, l)\) ，设
\end{enumerate}

\[
f _ {i k l}: \mathrm{H} _ {i k} ^ {\prime} \times \mathrm{H} _ {k l} ^ {\prime \prime} \rightarrow \mathrm{H} _ {i l}
\]

是一个映射。那么，若 \(M' = (m_{ik}')_{(i,k) \in I \times K}\) , \(M'' = (m_{kl}'')_{(k,l) \in K \times L}\) 是满足 \(m_{ik}' \in H_{ik}'\) 与 \(m_{kl}'' \in H_{kl}''\) （对所有 \(i, k, l\) ）的矩阵，我们就可以定义乘积 \(M'M''\) （经由诸 \(f_{ikl}\) ）。我们把写出并证明类似于 (4)、(5) 与 (6) 的公式的工作留给读者。

\subsection{环上的矩阵}

数学中最重要的矩阵是环 A 上的矩阵。对应于指标集 I, K 的 A 上矩阵的集合 \(A^{I \times K}\) 于是典范地具有一个 (A, A)-双模结构（§1，第 14 号）。

对每个有序对 \((i,k)\in\mathbf{I}\times\mathbf{K}\) ，设 \(E_{ik}\) 是矩阵 \((a_{jl})\) ，它满足 \(a_{ik}=1\) 且 \(a_{jl}=0\) （对 \((j,l)\neq(i,k)\) ）；诸 \(E_{ik}\) 称为矩阵集合 \(\mathbf{A}^{\mathbf{I}\times\mathbf{K}}\) 中的矩阵单位；若 I 与 K 有限，它们构成该集合关于其左或右 A-模结构的典范基（§1，第 11 号）。显然

\[
{ } ^ { t } E _ { i k } = E _ { k i } .
\]

除非另有说明，A 上两个矩阵（假定有定义）的乘积 \(M'M''\) 将总是理解为相对于 A 中乘法 \((x,y)\mapsto xy\) 的乘积（或者说，“在 A 中计算”）。于是我们有（第 2 号）结合律与分配律公式

\[
(X Y) Z = X (Y Z)\tag{7}
\]

\[
\left\{ \begin{array}{l} X (Y + Z) = X Y + X Z \\ (X + Y) Z = X Z + Y Z \end{array} \right.\tag{8}
\]

对于 A 上的三个矩阵 X, Y, Z，只要这些公式中出现的和与积有定义。

特别是，如果 \(E_{ik}\) （相应地 \(E_{kl}'\) , \(E_{il}''\) ）分别是 \(\mathbf{A}^{\mathbf{I} \times \mathbf{K}}\) （相应地 \(\mathbf{A}^{\mathbf{K} \times \mathbf{L}}\) , \(\mathbf{A}^{\mathbf{I} \times \mathbf{L}}\) ）中的矩阵单位，且 \(\mathbf{I} = [1, p]\) , \(\mathbf{K} = [1, q]\) , \(\mathbf{L} = [1, r]\) ，我们得到公式

\[
\left\{ \begin{array}{l l} E _ {i k} E _ {j l} ^ {\prime} = 0 & \text { if } k \neq j \\ E _ {i k} E _ {k l} ^ {\prime} = E _ {i l} ^ {\prime \prime}. \end{array} \right.\tag{9}
\]

设 \(\mathbf{A}^0\) 是 \(\mathbf{A}\) 的反环，并用 \(a * b (= ba)\) 表示 \(a\) 与 \(b\) 在 \(\mathbf{A}^0\) 中的乘积；那么，对两个矩阵 \(X, Y\) （ \(\mathbf{A}\) 上的），其乘积有定义，

\[
{ } ^ { t } ( X Y ) = { } ^ { t } Y * { } ^ { t } X\tag{10}
\]

其中右边的 \(^{t}Y\) 与 \(^{t}X\) 被看作元素属于 \(A^{0}\) 的矩阵；当 A 交换时，则

\[
{ } ^ { t } ( X Y ) = { } ^ { t } Y . { } ^ { t } X\tag{11}
\]

命题 1. 设 A、B 是两个环，且 \(M = (m_{ik})_{(i,k) \in I \times K}\) 和

\[
M ^ {\prime} = \left(m _ {i k} ^ {\prime}\right) _ {(i, k) \in \mathbf {I} \times \mathbf {K}}
\]

两个具有有限指标集的矩阵，取值于一个 (A, B)-双模 G 上。假设对每个单行、元素属于 A 的矩阵单位 \(L = (a_{i})_{i \in I}\) 以及每个单列、元素属于 B 的矩阵单位 \(C = (b_{k})_{k \in K}\) ，都有 L.M.C = L.M'.C（乘积经由 (A, B)-模 G 的外运算计算）；则 \(M = M'\) 。

若把 \(L\) 取为矩阵单位 \((a_s)\) ，其中 \(a_i = 1\) , \(a_s = 0\) （对 \(s \neq i\) ），并把 \(C\) 取为矩阵单位 \((b_t)\) ，其中 \(b_k = 1\) , \(b_t = 0\) （对 \(t \neq k\) ），则乘积 \(L.M.C\) 与 \(L.M'.C\) 都是单元素矩阵，其元素分别等于 \(m_{ik}\) 与 \(m'_{ik}\) 。

设 A、B 是两个环， \(\sigma: A \to B\) 是一个同态。

对每个矩阵 \(M = (m_{\iota\kappa})\) （ \(\mathbf{A}\) 上的），我们将用 \(\sigma(M)\) 表示矩阵 \((\sigma(m_{\iota\kappa}))\) （ \(\mathbf{B}\) 上的）；显然 \(\sigma(aM) = \sigma(a)\sigma(M)\) , \(\sigma(Ma) = \sigma(M)\sigma(a)\) （对 \(a \in \mathbf{A}\) ），同时 \(\sigma(^t M) = ^t(\sigma(M))\) ，并且

\[
\left\{ \begin{array}{c} \sigma (M + M ^ {\prime}) = \sigma (M) + \sigma (M ^ {\prime}) \\ \sigma (M M ^ {\prime}) = \sigma (M) \sigma (M ^ {\prime}) \end{array} \right.\tag{12}
\]

当所考虑的运算有定义时，(12) 左边与右边的乘积分别在 A 与 B 中计算。当 \(\sigma\) 记作 \(x \mapsto x^{\sigma}\) 时，我们用 \(M^{\sigma}\) 代替 \(\sigma(M)\) 。

特别地，考虑 A 的一个反自同态 \(\sigma\) ，即 A 到反环 \(\mathbf{A}^0\) 的一个同态，或 A 到自身的一个映射，使得

\[
\sigma (a + a ^ {\prime}) = \sigma (a) + \sigma (a ^ {\prime}), \quad \sigma (a a ^ {\prime}) = \sigma (a ^ {\prime}) \sigma (a)
\]

对所有 \(a, a'\) （在 \(\mathbf{A}\) 中）；那么，对两个矩阵 \(M, M'\) （ \(\mathbf{A}\) 上的），其乘积 \(MM'\) 有定义时，

\[
\sigma (M M ^ {\prime}) = ^ {t} (\sigma (^ {t} M ^ {\prime}). \sigma (^ {t} M))\tag{13}
\]

其中两边的乘积均在A中计算；这由(10)和(12)直接得出。

\subsection{矩阵与线性映射}

设 A 是一个环，E 是一个具有基 \((e_i)_{i\in I}\) 的（右或左）A-模。对每个元素 \(x\in E\) ， \(x\) 关于基 \((e_i)\) 的矩阵，记作 \(M(x)\) 或 \(\mathbf{x}\) （当不致引起混淆时有时也简记为 \(x\) ），是由分量 \(x_i\) ( \(i\in I\) )（即 \(x\) 关于 \((e_i)\) 的分量）组成的列矩阵（ \(\S 1\) ，第 11 号）；在计算中，为了记住指标 \(i\) 是行指标，有时方便给它添上一个只取一个值的列指标，并把矩阵 \(M(x)\) 写成 \((x_{i0})\) 。

我们现在考虑两个（左或右）A-模 E 与 F，分别带有基 \((e_{i})_{i\in I}\) 与 \((f_{k})_{k\in K}\) ；令 \((f_{k}^{*})\) 为对应于 \((f_{k})\) 的坐标形式族。对 E 到 F 的线性映射 u，我们将在以下每种情形定义 u 关于基 \((e_{i})\) , \((f_{k})\) 的矩阵：

\begin{enumerate}
\def\labelenumi{(\Alph{enumi})}
\setcounter{enumi}{3}
\item
  E和F是右A模，u是A线性的。
\item
  E和F是左A模，u是A线性的。
\end{enumerate}

在下文中，我们将字母(D)（相应地(G)）附加到适用于右（相应地左）模的公式上。

定义 3. 在上述两种情形的每一种中， \(u\) 关于基 \((e_i), (f_k)\) 的矩阵是矩阵 \(M(u) = (u_{ki})_{(k,i) \in K \times I}\) ，使得

\[
u _ {k i} = f _ {k} ^ {*} (u (e _ {i}))\tag{14}
\]

分别写为

\[
u _ {k i} = \langle f _ {k} ^ {*}, u (e _ {i}) \rangle\tag{14 D}
\]

\[
u _ {k i} = \langle u (e _ {i}), f _ {k} ^ {*} \rangle .\tag{14 G}
\]

于是 \(M(u)\) 的指标为 i 的列等于 \(M(u(e_{i}))\) 。

显然，若 u, v 是 E 到 F 的两个线性映射，而 \(M(u)\) , \(M(v)\) 是它们关于同一组基的矩阵，则

\[
M (u + v) = M (u) + M (v)\tag{15}
\]

和

\[
M (\gamma u) = \gamma M (u)\tag{16}
\]

对每个元素 \(\gamma\) （属于 A 的中心 \(\Gamma\) ）公式 (16) 成立。换言之，一旦基 \((e_i)\) , \((f_k)\) 固定，映射 \(u \mapsto M(u)\) 就是一个 \(\Gamma\) -模同构，从 \(\operatorname{Hom}_{\mathbf{A}}(\mathbf{E}, \mathbf{F})\) 映到集合 \(\mathbf{A}^{\mathbf{K} \times \mathbf{I}}\) 的一个子集上，该子集等于 \(\mathbf{A}^{\mathbf{K} \times \mathbf{I}}\) （若 \(\mathbf{K}\) 有限）。

命题 2. 设 I 与 K 有限。对每个元素 \(x \in \mathbf{E}\) ，矩阵 \(M(u(x))\) 关于基 \((f_k)\) 由以下公式给出

\[
M (u (x)) = M (u). M (x)\tag{17 D}
\]

\[
{ } ^ { t } M ( u ( x ) ) = { } ^ { t } M ( x ) . { } ^ { t } M ( u ) .\tag{17 G}
\]

我们例如验证 (17 G)。设 \(x = \sum_{i} x_{i0} e_i\) , \(u(x) = \sum_{k} y_{k0} f_k\) ，其中 \(x_{i0} \in \mathbf{A}\) , \(y_{k0} \in \mathbf{A}\) ；于是 \(u(x) = u\left(\sum_{i} x_{i0} e_i\right) = \sum_{i} x_{i0} u(e_i) = \sum_{i, k} x_{i0} u_{ki} f_k\) ；由此 \(y_{k0} = \sum_{i} x_{i0} u_{ki}\) 。为把两个指标 \(i\) 并排放在一处，我们考虑转置矩阵 \(^t M(x) = (x'_{0i})\) ，其中 \(x'_{0i} = x_{i0}\) ，以及

\[
{ } ^ { t } M ( u ) = \left( u _ { i k } ^ { \prime } \right) ,
\]

其中 \(u_{ik}' = u_{ki}\) ；于是 \(y_{k0} = \sum_{i} x_{0i}' u_{ik}'\) ，而右边是单行矩阵中指标为 \(k\) 的元素（该矩阵为 \(^t M(x).^t M(u)\) ），由此得 (17 G)。

当 A 交换时，(17 G) 由第 2 号的公式 (4) 化为 (17 D)。

推论。设 E, F, G 是环 A 上的三个右（相应地左）模， \((e_{i})_{i\in I}\) , \((f_{k})_{k\in K}\) , \((g_{l})_{l\in L}\) 分别是 E、F、G 的有限基， \(u:E\to F\) , \(v:F\to G\) 是两个线性映射， \(M(u)\) 是 u 关于基 \((e_{i})\) , \((f_{k})\) 的矩阵， \(M(v)\) 是 v 关于基 \((f_{k})\) , \((g_{l})\) 的矩阵，而 \(M(v\circ u)\) 是 \(v\circ u\) 关于基 \((e_{i})\) , \((g_{l})\) 的矩阵；则

\[
M (v \circ u) = M (v) M (u)\tag{18 D}
\]

\[
{ } ^ { t } M ( v \circ u ) = { } ^ { t } M ( u ) { } ^ { t } M ( v ) .\tag{18 G}
\]

我们以(18 G)为例来证明。对所有 \(x \in \mathbf{E}\) ，由(17 G)：

\(^{t}M(x).^{t}M(v\circ u)=^{t}M(v(u(x)))=^{t}M(u(x)).^{t}M(v)=^{t}M(x).^{t}M(u).^{t}M(v)\) 根据结合性；然后推论由第3号命题1得出，因为矩阵 \(^{t}M(x)\) 只有一行，是任意的。

注 (1)。公式 (17 D) 可视为 (18 D) 的特例。因为对每个 \(x \in \mathbf{E}\) ，典范地对应着线性映射 \(\theta_x: \mathbf{A}_d \to \mathbf{E}\) ，它把每个 \(\alpha \in \mathbf{A}\) 映为 \(x\alpha\) （§ 2，第 1 号）。立即可见，矩阵 \(M(\theta_x)\) （关于 \(\mathbf{A}_d\) 的基 1 以及基 \((e_i)\) （ \(\mathbf{E}\) 的） ）正是矩阵 \(M(x)\) ；类似地 \(M(\theta_{u(x)}) = M(u(x))\) ，因此公式 (17 D) 可视为关系式

\[
\theta_ {u (x)} = u \circ \theta_ {x}.
\]

命题 3. 设 E、F 是两个右（相应地左）A-模，且 \((e_i)_{i\in I}, (f_k)_{k\in K}\) 分别是 E 与 F 的有限基。对 E 到 F 的每个线性映射 \(u\) ，设 \(M(u)\) 是 \(u\) 关于基 \((e_i)\) 与 \((f_k)\) 的矩阵。则 \(^t u: F^* \to E^*\) 关于对偶基 \((f_k^*)\) 与 \((e_i^*)\) 的矩阵等于 \(^t M(u)\) 。

E 典范地等同于它的双对偶 \(E^{**}\) ，且 \((e_{i})\) 等同于 \((e_{i}^{*})\) 的对偶基；于是（例如设 E 与 F 是右模）

\[
\langle^ {t} u (f _ {k} ^ {*}), e _ {i} \rangle = \langle f _ {k} ^ {*}, u (e _ {i}) \rangle ,
\]

由此得出该命题。

注记。(2) 设 E 与 F 是两个左 A-模，分别具有基 \((e_{i})_{i\in\mathbf{I}}\) 与 \((f_{k})_{k\in\mathbf{K}}\) 。对每个 A-线性映射 \(u:E\to F\) ，由 (14 G) 有 \(u(e_{i})=\sum_{k}u_{ki}f_{k}\) ；这些关系也可这样解释：元素属于 F 的列矩阵 \((u(e_{i})_{i\in\mathbf{I}})\) 等于乘积 \({}^{t}\mathbf{M}(u).(f_{k})\) ，其中 \((f_{k})_{k\in\mathbf{K}}\) 被看作元素属于 F 的列矩阵，且乘积是对映射 \(A\times F\to F\) （它定义 A-模 F 上的作用律，第 2 号）计算的。

\begin{enumerate}
\def\labelenumi{(\arabic{enumi})}
\setcounter{enumi}{2}
\tightlist
\item
设 A、B 是两个交换环， \(\sigma: A \to B\) 是一个环同态。在命题 3 的记号下， \((e_i \otimes 1)\) 与 \((f_k \otimes 1)\) 分别是 \(E_{(B)} = E \otimes_A B\) 与 \(F_{(B)} = F \otimes_A B\) 的基（ \(\S 5\) ，第 1 号，命题 4）；此外，若 \((e_i^*)\) 与 \((f_k^*)\) 分别是 \((e_i)\) 与 \((f_k)\) 的对偶基，则 \((e_i^* \otimes 1)\) 与 \((f_k^* \otimes 1)\) 分别是 \((e_i \otimes 1)\) 与 \((f_k \otimes 1)\) 的对偶基（ \(\S 5\) ，第 4 号）。对每个 A-线性映射 \(u: E \to F\) ，设 \(M(u)\) 与 \(M(u \otimes 1)\) 分别是 \(u\) 关于 \((e_i)\) 与 \((f_k)\) 的矩阵以及 B-线性映射 \(u \otimes 1\) 关于 \((e_i \otimes 1)\) 与 \((f_k \otimes 1)\) 的矩阵。由 \(\S 5\) 第 4 号公式 (20) 可知
\end{enumerate}

\[
M (u \otimes 1) = \sigma (M (u)).
\]

考虑一个含有限个未知量的、由有限个右纯量线性方程组成的方程组

\[
\sum_ {i \in I} a _ {k i} x _ {i} = b _ {k} \quad (k \in K)\tag{19}
\]

其中 \(a_{ki}\) , \(x_{i}\) , \(b_{k}\) 属于 A。

设 \((e_i)_{i\in I}, (f_k)_{k\in K}\) 是 \(\mathbf{E} = \mathbf{A}_d^{\mathrm{I}}\) 与 \(\mathbf{F} = \mathbf{A}_d^{\mathrm{K}}\) 的典范基；方程组 (19) 等价于方程 \(u(x) = b\) ，其中 \(x = \sum_{i} e_i x_i\) , \(b = \sum_{k} f_k b_k\) ，而 \(u: \mathbf{E} \to \mathbf{F}\) 是使矩阵 \(M(u)\) （关于基 \((e_i)\) 与 \((f_k)\) ）等于 \(\mathbf{A} = (a_{ki})_{(k,i)\in K\times I}\) 的线性映射。这个矩阵称为线性方程组 (19) 的矩阵。回顾（§2，第 8 号，注记 2 与 3），若记 \(c_i = \sum_{k} f_k a_{ki}\) ，则方程组 (19) 等价于唯一的向量方程

\[
\sum_ {i} c _ {i} x _ {i} = b,\tag{20}
\]

且由于 \(c_{i}\) 是矩阵 A 中指标为 i 的列，可见说方程组 (19) 有解就等于说单列矩阵 \(b = (b_{k0})\) 是矩阵 A 的各列的线性组合。

我们把对左线性方程组作相应定义和注记的任务留给读者。

\subsection{分块乘积}

第 4 号的定义可以如下推广。设 E 是一个（右或左）A-模，是一个子模族 \((\mathbf{E}_i)_{i\in I}\) 的直和。对所有 \(x\in \mathbf{E}\) ，设 \(x = \sum_{i\in I}x_i\) ，其中 \(x_{i}\in \mathbf{E}_{i}\) （对所有 \(i\in I\) ）；我们称元素属于 E 的列矩阵 \(M(x) = (x_{i})_{i\in I}\) 为 \(x\) 关于 E 的直和分解 \((\mathbf{E}_i)_{i\in I}\) 的矩阵。

设 F 是另一个 A-模（E 与 F 同为右 A-模或同为左 A-模），并假设 F 是子模族 \((\mathbf{F}_{k})_{k\in\mathbb{K}}\) 的直和。对所有 \(u\in\mathrm{Hom}(\mathbf{E},\mathbf{F})\) 与所有 \(x_{i}\in\mathbf{E}_{i}\) ，设 \(u(x_{i})=\sum_{k}u_{ki}(x_{i})\) ，其中 \(u_{ki}(x_{i})\in\mathbf{F}_{k}\) （对所有 \(k\in K\) ）；于是 \(u_{ki}\in\mathrm{Hom}(\mathbf{E}_{i},\mathbf{F}_{k})\) ；我们称矩阵 \(M(u)=(u_{ki})_{(k,i)\in K\times I}\) —— (K, I) 型、元素属于集合 H（即诸 \(\mathrm{Hom}(\mathbf{E}_{i},\mathbf{F}_{k})\) 之和）—— 为 u 关于 E 与 F 的直和分解 \((\mathbf{E}_{i})\) 与 \((\mathbf{F}_{k})\) 的矩阵。

有了这些定义，显然若 \(u, v\) 是 \(E\) 到 \(F\) 的两个 A-线性映射，则对关于相同直和分解的矩阵，有

\[
M (u + v) = M (u) + M (v), \quad M (\gamma u) = \gamma M (u)\tag{21}
\]

（对 A 的中心的每个元素 \(\gamma\) ）（第 2 号，注记）。

此外，诸 \(u_{ki}\) 的定义表明，若 K 有限，则我们可写出

\[
M (u (x)) = M (u). M (x)\tag{22}
\]

其中 \(M(u(x))\) 是 \(u(x)\) 关于分解 \((\mathbf{F}_k)\) 的矩阵，(22) 右边的乘积是对映射 \((t, z) \mapsto t(z)\) （从 \(\operatorname{Hom}(\mathbf{E}_i, \mathbf{F}_k) \times \mathbf{E}_i\) 到 \(\mathbf{F}_k\) ）计算的（第 2 号，注记）。

设 G 是第三个 A-模，是一个子模族 \((\mathbf{G}_{l})_{l\in L}\) 的直和，于是对每个 A-线性映射 \(v:F\to G\) 对应着一个矩阵 \(M(v)=(v_{lk})\) （关于分解 \((\mathbf{F}_{k})\) 与 \((\mathbf{G}_{l})\) ）。若 I、K 与 L 都有限，则

\[
M (v \circ u) = M (v). M (u)\tag{23}
\]

其中左端是矩阵 \((w_{li})\) —— \(w = v \circ u\) 关于分解 \((\mathbf{E}_i)\) 与 \((\mathbf{G}_l)\) 的矩阵——，且左端的乘积是对映射 \((t, s) \mapsto t \circ s\) （从 \(\operatorname{Hom}(\mathbf{F}_k, \mathbf{G}_l) \times \operatorname{Hom}(\mathbf{E}_i, \mathbf{F}_k)\) 到 \(\operatorname{Hom}(\mathbf{E}_i, \mathbf{G}_l)\) ）计算的（第 2 号，注记）。这正是 §1 第 8 号的公式 (32) 用矩阵写出的形式。

最后，若假设 I 与 K 有限，则 \(E^{*}\) （相应地 \(F^{*}\) ）典范地等同于诸模 \(E_{i}^{*}\) （相应地 \(F_{k}^{*}\) ）的直和（ \(\S\) 2，第 6 号，命题 10）。于是立即可验证， \(^{t}u\) 关于分解 \((\mathbf{F}_{k}^{*})\) 与 \((\mathbf{E}_{i}^{*})\) 的矩阵正是 \((^{t}u_{kt})_{(k,t)\in\mathbb{K}\times\mathbb{I}}\) 。

现在设 I 与 K 有限，并且进一步设每个 \(E_{i}\) （相应地 \(F_{k}\) ）都具有一个有限基。这等价于说 E（相应地 F）具有一个基 \((e_{r})_{r\in\mathbb{R}}\) （相应地 \((f_{s})_{s\in\mathbb{S}}\) ），且 R（相应地 S）具有一个划分 \((\mathbf{R}_{i})_{i\in\mathbb{I}}\) （相应地 \((\mathbf{S}_{k})_{k\in\mathbb{K}}\) ），使得对所有 \(i\in I\) （相应地 \(k\in K\) ）， \((e_{r})_{r\in\mathbb{R}_{i}}\) 是 \(E_{i}\) 的基（相应地 \((f_{s})_{s\in\mathbb{S}_{k}}\) 是 \(F_{k}\) 的基）。于是，若 \(X=M(u)\) 是 u 关于基 \((e_{r})_{r\in\mathbb{R}}\) 与 \((f_{s})_{s\in\mathbb{S}}\) 的矩阵，则矩阵 \(X_{ki}=M(u_{ki})\) （关于基 \((e_{r})_{r\in\mathbb{R}_{i}}\) 与 \((f_{s})_{s\in\mathbb{S}_{k}}\) ）正是 X 的这样一个子矩阵：它由删去指标为 \(s\notin S_{k}\) 的行与指标为 \(r\notin R_{i}\) 的列而得到。这样我们便定义了一个一一对应

\[
X \mapsto (X _ {k i}) _ {(k, i) \in \mathbf {K} \times \mathbf {I}}\tag{24}
\]

介于元素属于 A 的 (S, R) 型矩阵的集合与矩阵的矩阵 \((\mathbf{X}_{ki})_{(k,i)\in\mathbf{K}\times\mathbf{I}}\) （ \(K\times I\) 型，其中每个 \(X_{ki}\) 是 \((\mathbf{S}_{k},\mathbf{R}_{i})\) 型的、元素属于 A 的矩阵）的集合之间。进一步设 G 具有一个有限基 \((g_{t})_{t\in\mathbf{T}}\) ，且 \(\mathbf{T}=(\mathbf{T}_{l})_{l\in\mathbf{L}}\) 是 T 的一个划分，使得对每个 \(l\in L\) ， \((g_{t})_{t\in\mathbf{T}_{l}}\) 是 \(G_{l}\) 的基；令 \(Y=M(v)\) 是 v 关于基 \((f_{s})_{s\in\mathbf{S}}\) 与 \((g_{t})_{t\in\mathbf{T}}\) 的矩阵， \(Y_{lk}=M(v_{lk})\) 是 \(v_{lk}\) 关于基 \((f_{s})_{s\in\mathbf{S}_{k}}\) 与 \((g_{t})_{t\in\mathbf{T}_{l}}\) 的矩阵，Z=M(w) 是 w=v∘u 关于基 \((e_{r})_{r\in\mathbf{R}}\) 与 \((g_{t})_{t\in\mathbf{T}}\) 的矩阵，而 \(Z_{li}=M(w_{li})\) 是 \(w_{li}\) 关于基 \((e_{r})_{r\in\mathbf{R}_{i}}\) 与 \((g_{t})_{t\in\mathbf{T}_{l}}\) 的矩阵；于是由 (23)，Z=YX 的子矩阵 \(Z_{li}\) 由下式给出

\[
Z _ {l i} = \sum_ {k} Y _ {l k} X _ {k i}\tag{25}
\]

换言之，当所涉及的所有乘积都有定义时，一一对应 (24) 把乘积变为乘积（矩阵的矩阵的乘积按第 2 号注记的意义定义）；当 Z=YX 的子矩阵 \(Z_{li}\) 如此计算时，就称该乘积是“按分块”进行的。

这个名称来自如下事实：当 \(I = [1, p]\) 且 \(K = [1, q]\) 时，表示矩阵 X 的表格被设想成划分成一些“块”，构成一个“矩阵的阵列”

\[
\left( \begin{array}{c c c c} X _ {1 1} & X _ {1 2} & \ldots & X _ {1 p} \\ X _ {2 1} & X _ {2 2} & \ldots & X _ {2 p} \\ \cdot & \cdot & \cdot & \cdot \\ X _ {q 1} & X _ {q 2} & \ldots & X _ {q p} \end{array} \right)
\]

当 p 与 q 是足够小的特定整数而使此做法可行时，它被看作表示 X 的符号。

\subsection{半线性映射的矩阵}

设 A、B 是两个环， \(\sigma: \mathbf{A} \to \mathbf{B}\) 是 A 到 B 的一个同态，E 是具有基 \((e_i)_{i \in I}\) 的右（相应地左）A-模，F 是具有基 \((f_k)_{k\in K}\) 的右（相应地左）B-模。设 \(u:\mathbf{E}\to \mathbf{F}\) 是相对于 \(\sigma\) 的半线性映射，且 \(u(e_i) = \sum_{k\in K}f_ku_{ki}\left(\text{resp.} u(e_i) = \sum_{k\in K}u_{ki}f_k\right)\) ，于是诸 \(u_{ki}\) 是 B 的元素；按定义，矩阵 \(M(u) = (u_{ki})\) （ \(\mathbf{K}\times \mathbf{I}\) 型）也称为 \(u\) 关于基 \((e_i)\) 与 \((f_k)\) 的矩阵。用与第 4 号命题 2 相同的计算，立即可验证对所有 \(x\in \mathbf{E}\) （若 I 与 K 有限）

\[
M (u (x)) = M (u). \sigma (M (x))\tag{26 D}
\]

（相应地

\[
{ } ^ { t } M ( u ( x ) ) = \sigma ( { } ^ { t } M ( x ) ) . { } ^ { t } M ( u ) ) .\tag{26 G}
\]

设 C 是第三个环， \(\tau: \mathbf{B} \to \mathbf{C}\) 是一个同态，G 是具有基 \((g_l)_{l \in \mathbf{L}}\) 的右（相应地左）C-模，而 \(v\) 是 F 到 G 的、相对于 \(\tau\) 的半线性映射；若 \(M(v)\) 是 \(v\) 关于 \((f_k)\) 与 \((g_l)\) 的矩阵，而 \(M(v \circ u)\) 是 \(v \circ u\) 相对于 \((e_i)\) 与 \((g_l)\) 的矩阵，则当 I、K 与 L 有限时，

\[
M (v \circ u) = M (v). \tau (M (u))\tag{27 D}
\]

（相应地

\[
{ } ^ { t } M ( v \circ u ) = \tau ( { } ^ { t } M ( u ) ) . { } ^ { t } M ( v ) ) .\tag{27 G}
\]

例如为证明 (27 D)，注意对所有 \(x \in E\) ，由 (26 D) 有

\[
\begin{array}{r l} M (v \circ u) \cdot \tau (\sigma (M (x))) & = M (v (u (x))) \\ & = M (v) \cdot \tau (M (u (x))) = M (v) \cdot \tau (M (u)) \cdot \tau (\sigma (M (x))), \end{array}
\]

从而由第 3 号的命题 1 得出 (27 D)。

最后设 \(\sigma: \mathbf{A} \to \mathbf{B}\) 是一个同构；那么回顾 \(^t u: \mathbf{F}^* \to \mathbf{E}^*\) 是相对于 \(\sigma^{-1}\) 的半线性映射（§ 2，第 5 号）；当 I 与 K 有限时， \(^t u\) 关于对偶基 \((f_k^*)\) 与 \((e_i^*)\) 的矩阵由下式给出

\[
M (^ {t} u) = \sigma^ {- 1} (^ {t} M (u))\tag{28}
\]

因为按定义（例如设 E 与 F 是右模）， \(\langle^t u(f_k^*), e_i \rangle^\sigma = \langle f_k^*, u(e_i) \rangle\) 当 \(\sigma\) 记作 \(x \mapsto x^\sigma\) 时成立。

注记。设 A 是一个环， \(\sigma\) 是 A 的一个反自同态（第 3 号）；考虑以下两种情形：

（GD）E 是左 A-模，F 是右 A-模，u 是 E 到 F 的一个 Z-线性映射，使得 \(u(ax) = u(x)\sigma(a)\) （对 \(a \in A\) , \(x \in E\) ）；换言之，u 是相对于 \(\sigma\) 的、右 \(A^{0}\)-模 E 到右 A-模 F 中的半线性映射。

（DG）E 是右 A-模，F 是左 A-模，u 是 E 到 F 的一个 Z-线性映射，使得 \(u(xa) = \sigma(a)u(x)\) （对 \(a \in A\) , \(x \in E\) ）；换言之，u 是相对于 \(\sigma\) 的、左 \(A^{0}\)-模 E 到左 A-模 F 中的半线性映射。

在这两种情形下，u 关于 E 与 F 的基的矩阵 \(M(u)\) ，其元素都属于 A；若这些基有限，则对所有 \(x \in E\) ，分别有公式

\[
M (u (x)) = M (u). \sigma (M (x))\tag{17 GD}
\]

\[
{ } ^ { t } M ( u ( x ) ) = \sigma ( { } ^ { t } M ( x ) ) . { } ^ { t } M ( u ) ,\tag{17 DG}
\]

两边的乘积都在 A 中计算。这分别由 (26 D) 与 (26 G) 直接得出。

\subsection{方阵}

定义 4. 行与列的指标集相同的矩阵称为方阵。

具有 n 行 n 列的方阵称为 n 阶矩阵。

注记。应当指出，行与列的指标集具有相同基数但并不相同的矩阵，绝不可视为方阵；特别地，这样的两个矩阵在环上的乘积是没有定义的。

显然，A 上的、以有限集作为行与列指标集的方阵的加法与乘法，由于公式 (7)、(8) 与 (9)（第 3 号），在这些矩阵的集合上定义了一个环结构；矩阵 \((\delta_{ij})\) （其中 \(\delta_{ij}\) 是 Kronecker 指标（对 \(i \in I\) , \(j \in I\) ））是该环的单位元，记作 \(I_n\) 或 \(1_n\) （当 I 有 \(n\) 个元素时）。当 \(\mathbf{I} = [1, n]\) 时，如此定义的矩阵环简记为 \(\mathbf{M}_n(\mathbf{A})\) ； \(\mathbf{M}_n(\mathbf{A})\) 的可逆元素群记作 \(\mathbf{GL}_n(\mathbf{A})\) 或 \(\mathbf{GL}(n, \mathbf{A})\) 。

A 上的 n 阶方阵 \(U = (a_{ij})\) 为右（相应地左）可逆的必要且充分条件是：对 A 的元素的每个组 \((b_i)_{1 \leq i \leq n}\) ，由 n 个未知量的 n 个方程构成的方程组

\[
\sum_ {j = 1} ^ {n} a _ {i j} x _ {j} = b _ {i} \quad (1 \leqslant i \leqslant n)
\]

\[
\left(\text {resp.} \sum_ {j = 1} ^ {n} x _ {j} a _ {j i} = b _ {i}\right)
\]

都在 A 中有一个解 \((x_{i})\) 。

设 I 是有限指标集，A 是环，E 是具有基 \((e_{i})_{i\in I}\) 的右（相应地左）A-模。对 E 的每个自同态 u，矩阵 \(M(u)\) （u 关于两个都等同于 \((e_{i})\) 的基的）是一个方阵；更简单地说，它称为 u 关于基 \((e_{i})\) 的矩阵。

设 \(\mathbf{I} = [1, n]\) 。映射 \(u \mapsto M(u)\) （相应地 \(u \mapsto {}^t M(u)\) ）是环 \(\operatorname{End}_{\mathbf{A}}(\mathbf{E})\) 到 \(\mathbf{M}_n(\mathbf{A})\) （相应地到 \(\mathbf{M}_n(\mathbf{A})\) 的反环上，这由公式 (18 D)（相应地 (18 G)）得出）（第 4 号）的一个同构。环 \(\mathbf{M}_{n}(\mathbf{A})\) 的可逆元素称为可逆矩阵，它们在映射 \(u \mapsto M(u)\) （相应地 \(u \mapsto {}^{t}M(u)\) ）之下对应于 E 的自同构；因此群 \(\mathbf{GL}(n, \mathbf{A})\) 典范地等同于群 \(\mathbf{GL}(\mathbf{A}_{d}^{n})\) 。

若 \(u\) 是 \(\mathbf{E}\) 的一个自同构，则它的逆步 \(\check{u}\) 是左（相应地右）A-模 \(\mathbf{E}^*\) 的一个自同构，使得 \(\check{u} = (^{t}u)^{-1} = ^{t}(u^{-1})\) （§ 2，第 5 号，定义 6）；若 \(M(\check{u})\) 是 \(\check{u}\) 关于对偶基 \((e_{i}^{*})\) 的矩阵，则由命题 3（第 4 号）可得

\[
M (\check {u}) = (^ {t} M (u)) ^ {- 1} = ^ {t} M (u ^ {- 1}).\tag{29}
\]

因此，对每个可逆矩阵 X，有 \({}^{t}(X^{-1}) = ({}^{t}X)^{-1}\) ；这个矩阵也记作 \({}^{t}X^{-1}\) ，称为矩阵 X 的逆步矩阵。

设 \(\sigma\) 是环 A 的一个自同构；对每个半线性映射 \(u: \mathbf{E} \to \mathbf{E}\) （相对于 \(\sigma\) ），矩阵 \(M(u)\) （该映射关于 E 的基 \((e_i)\) 的）也是一个方阵。由 (27 D)（第 6 号）立即可得，若 \(u\) 是双射，则

\[
M (u ^ {- 1}) = (\sigma^ {- 1} (M (u))) ^ {- 1}.
\]

设 E 是一个 A-模，它是子模的有限族 \((\mathbf{E}_{i})_{i\in I}\) 的直和；对 E 的每个自同态 u，矩阵 \(M(u)=(u_{kt})\) （u 关于 E 的两个都等同于 \((\mathbf{E}_{i})\) 的分解（第 5 号）的）是一个线性映射的方阵。要使 \(u(\mathbf{E}_{i})\subset\mathbf{E}_{i}\) 对所有 \(i\in I\) 成立，必要且充分条件是 \(u_{kt}=0\) （对 \(k\neq i\) ）。当 \(\mathbf{I} = [1, n]\) 时，关系

\[
u (\mathrm{E} _ {i}) \subset \mathrm{E} _ {i} + \mathrm{E} _ {i + 1} + \dots + \mathrm{E} _ {n} \quad (1 \leqslant i \leqslant n)
\]

等价于关系 \(u_{ki} = 0\) （对 \(k < i\) ）。

方阵的例子。I. 对角矩阵。在一个方阵中，

\[
M = \left(m _ {\iota \kappa}\right) _ {(\iota, \kappa) \in I \times I},
\]

两个指标相等的元素称为对角元素，族 \((m_{\iota})_{\iota \in I}\) 称为 \(M\) 的对角线；环上的方阵 \(M = (m_{\iota \kappa})\) ，若其对角元素以外的元素全为零，则称为对角矩阵。对环 A 的元素的每个族 \((a_{\iota})_{\iota \in I}\) ，对角矩阵 \((m_{\iota \kappa})\) （满足 \(m_{\iota} = a_{\iota}\) （对所有 \(\iota \in I\) ））记作 \(\operatorname{diag}(a_{\iota})_{\iota \in I}\) （或 \(\operatorname{diag}(a_1, a_2, \ldots, a_n)\) ，当 \(I = [1, n]\) 时）。在集合 \(\mathbf{M}_n(A)\) （A 上 \(n\) 阶方阵的）中，单位矩阵 \(I_n\) 是对角矩阵，并且它的每个倍数 \(aI_n = I_na\) （该矩阵被纯量 \(a\) 乘）也是（即称为标量矩阵的对角矩阵，其所有对角元素都等于 \(a\) ）。

对 A 的元素的每个族 \((d_i)_{1\leqslant i\leqslant n}\) 以及每个矩阵 \(X = (x_{ij})\) （ \((n,q)\) 型（相应地 \((p,n))\) ），A 上的），记 \(D = \mathrm{diag}(d_i)\) ，则

\[
\left\{ \begin{array}{l} D X = (d _ {i} x _ {i j}) \\ X D = (x _ {i j} d _ {j}). \end{array} \right.\tag{30}
\]

特别地，对两个 n 阶对角矩阵，有

\[
\begin{array}{c} \operatorname{diag} (a _ {i}) + \operatorname{diag} (b _ {i}) = \operatorname{diag} (a _ {i} + b _ {i}) \\ \operatorname{diag} (a _ {i}). \operatorname{diag} (b _ {i}) = \operatorname{diag} (a _ {i} b _ {i}). \end{array}\tag{31}
\]

因此对角矩阵构成 \(\mathbf{M}_{n}(\mathbf{A})\) 的一个子环，它同构于乘积环 \(A^{n}\) ；标量矩阵构成一个同构于 A 的子环。

\begin{enumerate}
\def\labelenumi{\Roman{enumi}.}
\setcounter{enumi}{1}
\tightlist
\item
  置换矩阵；单项矩阵。设 \(\pi\) 是有限集 I 的任意一个置换，并设 \((e_i)_{i\in I}\) 是 A-模 \(\mathbf{E} = \mathbf{A}_d^{\mathbf{I}}\) 的典范基；存在且只存在一个 E 的自同态 \(u_{\pi}\) ，使得对所有 \(i\in I\) , \(u_{\pi}(e_i) = e_{\pi(i)}\) （§ 1，第 11 号，命题 17 的推论 3）。对所有 \(i\in I\) ，指标为 \(i\) 的列（在矩阵 \(M(u_{\pi})\) 中，关于基 \((e_i)\) ）除指标为 \(\pi(i)\) 的行中那个等于 1 的元素外，其余元素全为零。按一种习用说法， \(M(u_{\pi})\) 称为置换 \(\pi\) 的矩阵。立即可见，对 I 的任意两个置换 \(\sigma, \tau\) ，有 \(u_{\sigma \tau} = u_{\sigma} \circ u_{\tau}\) ，且对恒等置换 \(\varepsilon, u_{\varepsilon}\) 为恒等映射；因此映射 \(\pi \mapsto M(u_{\pi})\) 是对称群 \(\mathfrak{S}_{\mathbf{I}}\) 到置换矩阵群上的一个同构。
\end{enumerate}

置换矩阵的每一行与每一列只含一个元素 \(\neq0\) 。非零环 A 上的、具有此性质的有限方阵 R 称为单项矩阵；设 \(r_{i}\) 是 R 的指标为 i 的列中那个唯一的元素 \(\neq0\) ，并设 \(\pi(i)\) 为该元素所在行的指标；显然 \(\pi\) 是指标集 I 的一个置换，且 \(R=M(u_{\pi})D\) ，其中 \(D=\mathrm{diag}(r_{i})\) 。

\begin{enumerate}
\def\labelenumi{\Roman{enumi}.}
\setcounter{enumi}{2}
\tightlist
\item
  三角矩阵。在环 A 上 n 阶方阵的环 \(\mathbf{M}_{n}(\mathbf{A})\) 中，任意矩阵 \((a_{ij})\) ，若 \(a_{ij}=0\) （对 i\textgreater j（相应地 i\textless j）），则称为上（相应地下）三角矩阵；也说这样的矩阵在其对角线以下（相应地以上）只有零。立即可证，上（相应地下）三角矩阵构成 \(\mathbf{M}_{n}(\mathbf{A}),\mathbf{S}\cap\mathbf{T}\) 的一个子环 S（相应地 T），后者显然就是对角矩阵环。
\end{enumerate}

S（相应地 T）中对角元素可逆的矩阵所成的集合 S'（相应地 T'）是一个矩阵乘法群，称为上（相应地下）全三角群，这由 § 1 第 11 号注记 5 直接得出。S（相应地 T）中对角元素全等于 1 的矩阵所成的集合 \(S_{1}\) （相应地 \(T_{1}\) ）是上述群的一个子群，称为上（相应地下）严格三角群，且 S'（相应地 T'）中对角线为 \((d_{i})\) 的每个矩阵 M 都可写成 \(M = DM_{1} = M_{1}'D\) ，其中 \(D = \mathrm{diag}(d_{i})\) ，而 \(M_{1}\) 与 \(M_{1}'\) 是属于 \(S_{1}\) （相应地 \(T_{1}\) ）的矩阵。

\begin{enumerate}
\def\labelenumi{\Roman{enumi}.}
\setcounter{enumi}{3}
\tightlist
\item
  矩阵的对角矩阵与三角矩阵。设 \((\mathbf{I}_k)_{1 \leqslant k \leqslant p}\) 是有限集 I 的一个划分；环 A 上以 I 为指标集的每个方阵都可以写成一个矩阵的方阵的形式，它对应于行的指标集与列的指标集的同一划分 \((I_{k})\) （第 5 号）
\end{enumerate}

\[
\left( \begin{array}{c c c c} X _ {1 1} & X _ {1 2} & \ldots & X _ {1 p} \\ X _ {2 1} & X _ {2 2} & \ldots & X _ {2 p} \\ \cdot & \cdot & \cdot & \cdot \\ X _ {p 1} & X _ {p 2} & \ldots & X _ {p p} \end{array} \right)\tag{32}
\]

其中每个 \(X_{kk}\) 是以 \(I_{k}\) 为行和列指标集的方阵。

在此记号下，若所有矩阵 \(X_{ij}\) （满足 \(i \neq j\) （相应地 i \textgreater{} j，相应地 i \textless{} j）的）均为零，则称 (32) 为矩阵的对角矩阵（相应地矩阵的上三角矩阵、相应地矩阵的下三角矩阵）。其矩阵为矩阵的对角矩阵、相应地三角矩阵的自同态 u 的解释，前面已通过考虑相应的线性映射的矩阵 \(M(u)\) 而看到过。对 I 的给定划分 \((\mathbf{I}_{k})\) ，矩阵的下三角（相应地上三角、对角）矩阵构成矩阵环 \(A^{I \times I}\) 的一些子环。特别地，相对于

划分 \((\mathbf{I}_k)\) 的矩阵的对角矩阵环同构于乘积 \(\prod_{k=1}^{p} \operatorname{End}_{\mathbf{A}}(\mathbf{E}_k)\) 。
